%%%%%%%%%%%%%%%%%%%%%%%%%%%%%%%%%%%%%%%%%
% The Legrand Orange Book
% LaTeX Template
% Version 2.3 (8/8/17)
%
% This template has been downloaded from:
% http://www.LaTeXTemplates.com
%
% Original author:
% Mathias Legrand (legrand.mathias@gmail.com) with modifications by:
% Vel (vel@latextemplates.com)
%
% License:
% CC BY-NC-SA 3.0 (http://creativecommons.org/licenses/by-nc-sa/3.0/)
%
% Compiling this template:
% This template uses biber for its bibliography and makeindex for its index.
% When you first open the template, compile it from the command line with the 
% commands below to make sure your LaTeX distribution is configured correctly:
%
% 1) pdflatex main
% 2) makeindex main.idx -s StyleInd.ist
% 3) biber main
% 4) pdflatex main x 2
%
% After this, when you wish to update the bibliography/index use the appropriate
% command above and make sure to compile with pdflatex several times 
% afterwards to propagate your changes to the document.
%
% This template also uses a number of packages which may need to be
% updated to the newest versions for the template to compile. It is strongly
% recommended you update your LaTeX distribution if you have any
% compilation errors.
%
% Important note:
% Chapter heading images should have a 2:1 width:height ratio,
% e.g. 920px width and 460px height.
%
%%%%%%%%%%%%%%%%%%%%%%%%%%%%%%%%%%%%%%%%%

%----------------------------------------------------------------------------------------
%	PACKAGES AND OTHER DOCUMENT CONFIGURATIONS
%----------------------------------------------------------------------------------------

\documentclass[11pt,fleqn]{book} % Default font size and left-justified equations
\usepackage[UTF8]{ctex}

\usepackage{subfigure}
\usepackage{float}
\usepackage{tcolorbox}
\usepackage{lipsum}
\usepackage{listings}
\usepackage[]{algorithm2e}
\usepackage{blindtext}
\usepackage{amsmath,amssymb,amsfonts}
\usepackage{framed}
\usepackage{tabularx}
\usepackage{longtable}
\usepackage{supertabular}%长表格绘制
%\usepackage{enumerate}
%\usepackage{hyperref}

%\usepackage{changes}
%\definechangesauthor[name={ZYF}, color=orange]{zyf}

%-----------------------------------------------------------------------
% This part is utilized to generate color box.
% These codes are writed by Dr. Zhan
%----------------------------------------------------------------------
\usepackage{paralist}

\newtcolorbox{red}{
notitle,
colback=red!10
}

\newtcolorbox{blue}{
notitle,
colback=blue!10
}

\newtcolorbox{green}{
notitle,
colback=green!10
}


\newtcolorbox{gray}{
notitle,
colback=gray!10
}
%--------------------------------------------------------------

%-----------------------------------------------------
\usepackage{yhmath}
%\usepackage{algorithm}
%\usepackage{algorithmicx,algpseudocode}

\usepackage{listings}
\usepackage{color}
\definecolor{dkgreen}{rgb}{0,0.6,0}
\definecolor{gray}{rgb}{0.5,0.5,0.5}
\definecolor{mauve}{rgb}{0.58,0,0.82}
\lstset{frame=tb,
	language=Python,
	aboveskip=3mm,
	belowskip=3mm,
	showstringspaces=false,
	columns=flexible,
	basicstyle={\small\ttfamily},
	numbers=left,%设置行号位置none不显示行号
	%numberstyle=\tiny\courier, %设置行号大小
	numberstyle=\tiny\color{gray},
	keywordstyle=\color{blue},
	commentstyle=\color{dkgreen},
	stringstyle=\color{mauve},
	breaklines=true,
	breakatwhitespace=true,
	escapeinside=``,%逃逸字符(1左面的键)，用于显示中文例如在代码中`中文...`
	tabsize=4,
	extendedchars=false %解决代码跨页时，章节标题，页眉等汉字不显示的问题
}

%\usepackage{changes}
%\definechangesauthor[name={ZYF}, color=orange]{zyf}

\usepackage{pythonhighlight}

%\renewcommand{\algorithmicrequire}{\textbf{Input:}}  % Use Input in the format of Algorithm
%\renewcommand{\algorithmicensure}{\textbf{Output:}} % Use Output in the format of Algorithm








%----------------------------------------------------------------------------------------

\input{structure} % Insert the commands.tex file which contains the majority of the structure behind the template

\begin{document}

%----------------------------------------------------------------------------------------
%	TITLE PAGE
%----------------------------------------------------------------------------------------

\begingroup
\thispagestyle{empty}
\begin{tikzpicture}[remember picture,overlay]
	\node[inner sep=0pt] (background) at (current page.center) {\includegraphics[width=\paperwidth]{background}};
	\draw (current page.center) node [fill=ocre!30!white,fill opacity=0.6,text opacity=1,inner sep=1cm]{\Huge\centering\bfseries\sffamily\parbox[c][][t]{\paperwidth}{\centering 小红书帖子热度研究报告\\[15pt] % Book title
			{\Large 实验计划书}\\[20pt] % Subtitle
			{\huge }}}; % Author name
\end{tikzpicture}
\vfill
\endgroup
%----------------------------------------------------------------------------------------
%	COPYRIGHT PAGE
%----------------------------------------------------------------------------------------

\newpage
~\vfill
\thispagestyle{empty}

\noindent Copyright \copyright\ 2021 Choujun Zhan\\ % Copyright notice

\noindent \textsc{Published by Publisher}\\ % Publisher

\noindent \textsc{book-website.com}\\ % URL

%\noindent Licensed under the Creative Commons Attribution-NonCommercial 3.0 Unported License (the ``License''). You may not use this file except in compliance with the License. You may obtain a copy of the License at \url{http://creativecommons.org/licenses/by-nc/3.0}. Unless required by applicable law or agreed to in writing, software distributed under the License is distributed on an \textsc{``as is'' basis, without warranties or conditions of any kind}, either express or implied. See the License for the specific language governing permissions and limitations under the License.\\ % License information
\noindent \textit{First printing, } % Printing/edition date

%----------------------------------------------------------------------------------------
%	TABLE OF CONTENTS
%----------------------------------------------------------------------------------------

%\usechapterimagefalse % If you don't want to include a chapter image, use this to toggle images off - it can be enabled later with \usechapterimagetrue

\chapterimage{chapter_head_1.pdf} % Table of contents heading image

\pagestyle{empty} % No headers

\tableofcontents % Print the table of contents itself

\cleardoublepage % Forces the first chapter to start on an odd page so it's on the right

\pagestyle{fancy} % Print headers again

%----------------------------------------------------------------------------------------
%	PART
%----------------------------------------------------------------------------------------

%\input{parts/new_content_tmp}
%----------------------------------------------------------------------------------------
%	PART
%----------------------------------------------------------------------------------------
\part{Part I: 研究背景与文献综述}
%----------------------------------------------------------------------------------------
%	CHAPTER 1
%----------------------------------------------------------------------------------------

\chapterimage{chapter_head_2.pdf} % Chapter heading image
\chapter{研究背景与文献综述}




\section{研究背景与意义}
随着移动互联网的深度普及和社交媒体平台的蓬勃发展，内容创作与分享已成为数字经济时代信息传播与消费决策的核心环节。小红书（RED）作为中国领先的生活方式分享社区和消费决策入口，已从早期的海外购物攻略平台，演变为集图文、短视频、直播等多种形式于一体的综合性内容生态。截至2023年底，小红书月活跃用户数已突破**3亿**，其中**70\%**为90后年轻用户，社区内日均笔记发布量高达**数百万篇**，形成了一个庞大、活跃且极具商业价值的内容市场。在此背景下，品牌方、商家及个人创作者如何在海量信息中脱颖而出，通过优化内容策略以实现更高的曝光度、互动量与转化效果，已成为一个至关重要且富有挑战性的现实课题。

然而，小红书平台的内容分发机制相对复杂且不透明，一篇笔记能否成为“爆款”，往往受到内容质量、呈现形式、发布时机、标签策略、用户互动（如初始点赞、收藏、转发）以及平台算法推荐等多重因素的共同影响。例如，一条看似普通的日常分享笔记，可能因契合了某个热点话题、拥有吸引眼球的封面图、或在特定时间段发布，而在短时间内获得指数级传播，实现“从0到1”的突破；反之，精心策划的商业推广内容若未能准确把握平台调性与用户偏好，则可能陷入“自娱自乐”的窘境。这种不确定性不仅增加了内容创作者与营销者的试错成本，也使得广告投放的精准性与投资回报率（ROI）难以稳定评估。

本研究的意义在于，基于小红书公开的笔记数据，系统地量化分析影响笔记热度的关键因素（如内容特征、视觉元素、时间地点、互动指标、话题标签等），并构建科学的数据分析模型，旨在揭示内容热度生成的内在逻辑与规律。通过实证研究，我们期望为内容创作者、品牌营销人员及广告代理商提供一套具有可操作性的内容优化策略框架，帮助其在内容策划、形式设计、发布规划及互动运营等环节做出更科学、更有效的决策，从而在竞争激烈的“种草”生态中提升内容的吸引力、传播力与商业转化力，最终实现广告效果的最大化。此外，本研究亦有助于深化对社交媒体内容传播机制的理解，为平台算法的优化与社区生态的健康发展提供来自数据侧的洞察。

\section{国内外研究现状}
\subsection{社交媒体内容传播与影响力研究综述}
社交媒体内容传播与影响力的研究已成为传播学、市场营销学及数据科学交叉领域的热点。其核心在于探究哪些因素驱动了内容的广泛传播（如转发、点赞、评论）以及用户参与（Engagement）。根据研究视角的差异，现有文献大致可分为三类。首先，基于内容特征的研究，关注文本情感、主题类别、语言风格、标题吸引力、信息实用性等对传播效果的影响。例如，研究发现，包含积极情绪或激起好奇心的标题往往能获得更高的点击率；而提供实用价值（如教程、攻略）的“干货”内容则更容易引发用户的收藏行为。其次，基于视觉与形式特征的研究，强调图片质量、视频时长、封面设计、排版美观度等多媒体元素的关键作用。例如，有研究表明，高质量的原创图片、统一和谐的视觉风格能显著提升用户的停留时间和互动意愿。再者，基于上下文与网络结构的研究，则聚焦于发布时间、发布者影响力（KOL/KOC）、社交网络关系、话题热度（Hashtag）以及平台算法推荐机制等外部环境因素。例如，识别用户活跃的高峰时段进行发布，或巧妙借势热点话题，已被证明是提升初始曝光量的有效策略。

近年来，随着数据获取与分析技术的进步，相关研究方法也在不断演进。传统的问卷调查与案例分析法正逐渐与大规模数据挖掘、自然语言处理（NLP）和计算机视觉（CV）技术相结合。例如，研究者利用情感分析模型量化笔记文本的情感倾向，利用图像识别技术分析封面图的元素构成（如人脸、产品、场景），进而建立其与互动数据的关联模型。同时，机器学习方法，特别是预测模型（如回归模型、分类模型）和因果推断方法（如自然实验、双重差分法），被越来越多地应用于量化单个或多个因素对内容传播效果的边际影响，并试图从相关性中识别出因果关系，为策略制定提供更坚实的证据。

\subsection{多模态社交媒体数据分析方法}
\subsubsection{传统统计与计量方法}
在社交媒体数据分析的早期阶段，传统统计与计量经济学方法扮演了基础角色。这些方法主要用于检验特定因素与内容表现（如点赞数、转发数）之间是否存在显著的相关关系。多元线性回归（Multiple Linear Regression）是最常用的模型之一，它将内容热度（因变量）表示为一系列可量化特征（自变量，如文本长度、图片数量、发布小时等）的线性组合，通过估计系数来评估各因素的贡献度。为了处理计数数据（如点赞数）的离散性和过度分散问题，泊松回归（Poisson Regression）或负二项回归（Negative Binomial Regression）也常被采用。

然而，社交媒体数据通常具有复杂的特征：因变量（如热度指标）分布高度偏斜（存在大量零值和少数极端的“爆款”）；自变量之间可能存在多重共线性（例如，高质量内容可能同时具备优质图文和合适的发布时间）；并且存在大量的潜在混淆变量。为了解决这些问题，研究者引入了更复杂的模型。例如，分位数回归（Quantile Regression）可用于研究不同热度水平下（如普通笔记 vs. 爆款笔记）影响因素的不同作用；而面板数据模型（Panel Data Model）则能利用同一创作者或同一主题的多次发布数据，控制个体或主题的固定效应，从而更准确地识别特定因素的作用。尽管这些方法在控制变量和估计参数方面非常严谨，但它们通常假设变量间的关系是线性的或可预定义形式的，对于捕捉文本、图像等高维复杂特征内部的非线性模式和交互作用能力有限。
\subsubsection{机器学习与深度学习方法}
机器学习和深度学习技术的发展，为分析和预测社交媒体内容热度提供了更强大、更灵活的工具。这些方法能够自动从原始数据（如原始文本、图像像素）中学习有效的特征表示，并建模复杂的非线性关系。

在特征工程方面，自然语言处理技术如词袋模型（Bag-of-Words）、TF-IDF以及词嵌入（Word Embedding，如Word2Vec, GloVe）可以将非结构化的文本转化为结构化特征，用于后续的回归或分类任务。对于图像，除了手动提取的颜色、纹理、物体检测等特征外，卷积神经网络（CNN）预训练模型（如ResNet, VGG）可以提取深层、抽象的视觉特征向量。

在预测建模方面，集成学习方法如随机森林（Random Forest）和梯度提升决策树（XGBoost, LightGBM）因其出色的性能和可解释性（通过特征重要性排序）而被广泛应用。这些模型能够处理高维特征，自动进行特征选择，并捕捉特征间的复杂交互作用，常用于构建内容热度的回归预测或分类（如预测笔记是否会成为热门）模型。

更进一步，深度学习模型特别适合处理端到端（End-to-End）的多模态学习任务。例如，可以构建融合文本（通过LSTM或Transformer编码器）和图像（通过CNN编码器）特征的神经网络模型，共同预测内容的互动量。图神经网络（GNN）则可用于建模用户-内容交互网络或话题传播网络，以纳入社交关系信息进行更精准的预测。例如，有研究尝试利用用户历史行为序列，通过循环神经网络（RNN）来个性化预测其对某条内容的感兴趣程度。这些先进方法不仅能提升预测的准确性，还能帮助研究者发现传统方法难以识别的深层模式和影响因素组合。

\section{内容热度预测与广告效果优化方法}
\subsubsection{传统预测与优化模型}
在内容热度预测的传统模型中，除了上述的统计回归模型，时间序列预测方法（如ARIMA模型）有时也被用于分析某个账号或话题的热度随时间变化的趋势，但这种方法更适用于宏观趋势分析，而非对单篇内容的微观预测。在广告效果优化方面，传统的A/B测试（A/B Testing）是实践中的黄金标准。通过将目标受众随机分为两组，分别展示不同版本的内容（如不同标题、不同封面图），然后比较两组在关键指标（如点击率、转化率）上的差异，可以科学地评估单一变量的影响。然而，A/B测试通常一次只能测试少数几个变量，且需要足够的流量和较长的测试周期，对于需要快速迭代、多因素同时优化的场景效率较低。

响应面方法（Response Surface Methodology, RSM）等实验设计方法可以用于探索多个因素（如标题长度、图片风格、发布时间）及其交互作用对结果指标的最优组合。但这类方法通常假设响应面是连续的且可用低阶多项式近似，在面对社交媒体内容这种离散、高维且存在大量噪声的优化问题时，其应用受到一定限制。
\subsubsection{机器学习驱动的内容策略优化}
机器学习方法为实现智能化、自动化的内容策略优化提供了新的路径。一方面，预测模型本身可以作为优化工具：通过构建一个能够相对准确预测内容热度的模型（即“代理模型”或“元模型”），创作者或营销者可以在内容发布前，模拟调整不同的内容属性（如尝试不同的标题关键词、选择不同的发布时间段），利用模型预测哪种组合可能获得更好的效果，从而指导决策。这本质上是一种基于模型的“假设分析”（What-if Analysis）。

另一方面，强化学习（Reinforcement Learning）为在动态环境中持续优化内容策略提供了框架。智能体（如一个自动发布系统）可以通过与环境的交互（发布内容 -> 观察用户反馈 -> 获得奖励），学习在不同情境下（如不同话题、不同受众时段）应采取的最优内容策略（如选择什么类型的内容、以何种形式呈现）。此外，因果森林（Causal Forest）等基于机器学习的因果推断方法，可以从观测数据中估计不同内容干预（如添加特定标签、使用特定滤镜）对个体笔记或特定用户群体的异质性处理效应，从而帮助实现更精细化的个性化内容优化。

最终，结合多模态数据分析、热度预测模型与因果推断/优化算法，本研究旨在构建一个系统性的分析框架，以科学解码小红书内容热度的生成密码，为提升广告与营销效果提供数据驱动的解决方案。

%----------------------------------------------------------------------------------------
%	PART
%----------------------------------------------------------------------------------------
\part{Part II: 小红书数据获取与数据处理}

%----------------------------------------------------------------------------------------
%	CHAPTER 1
%----------------------------------------------------------------------------------------

\chapterimage{chapter_head_2.pdf} % Chapter heading image

\chapter{数据获取与来源}

\section{数据来源}
小红书内容传播效果的研究高度依赖于高质量、大规模的真实用户笔记数据。为系统分析影响笔记互动表现（如点赞、评论、收藏等）的关键因素，本研究以小红书平台公开发布的用户笔记为主要研究对象，通过自动化数据采集方式构建研究所需的数据集。
在数据采集过程中，严格遵循相关法律法规及平台使用规范，仅采集用户公开可见的信息，并对涉及用户身份的数据进行脱敏处理，确保数据仅用于学术研究目的。
\subsection{主要数据源：小红书平台}
本研究的主要数据来源为小红书平台的公开笔记信息，数据采集覆盖家电与美妆两个典型内容垂直领域。通过对不同领域笔记内容及互动数据的对比分析，探究内容形式、文本特征及用户影响力等因素对笔记传播效果的影响。

数据采集时间范围为2023年1月至2024年8月，共获取约 1,338,527 条有效笔记数据，具有较大的样本规模和较强的统计分析价值。

采集的核心数据字段主要包括以下几类：
\begin{itemize}
	\item 笔记元数据：包括笔记唯一标识（ID）、笔记标题、正文内容、话题标签（关键词）、笔记发布时间及发布地点。
	\item 互动数据：包括笔记的点赞数、评论数和收藏数，用于衡量笔记的传播效果与用户参与度。
	\item 内容形式特征：包括笔记是否包含图片、是否包含视频及其链接信息，用于分析多媒体内容对传播效果的影响。
	\item 用户信息：包括博主名称及其粉丝数量，用于刻画发布者影响力水平，并作为后续分析中的控制变量。
\end{itemize}

% 插入单张图片
\begin{figure}[H]
    \centering 
    \includegraphics[width=0.8\textwidth]{note_page_fields.pdf}
    \caption{小红书笔记页面与采集字段示意 }
    \label{fig:note_page_fields}
\end{figure}

\begin{figure}[H] 
    \centering 
    \includegraphics[width=0.8\textwidth]{user_page_fields.pdf}
    \caption{用户页面采集字段示意 }
    \label{fig:user_page_fields}
\end{figure}

\subsection{数据领域划分方式}
为保证不同内容领域样本的准确性与可区分性，本研究在数据采集阶段即对内容领域进行前置区分，采用不同访问入口 URL分别获取家电与美妆领域的笔记数据。

其中，家电领域的数据通过平台搜索入口进行采集，基于以“家电”为核心关键词构建的搜索结果页面获取笔记列表；美妆领域的数据则通过平台官方内容频道入口进行采集，直接访问美妆频道对应的内容流页面获取笔记列表。由于不同入口页面在平台内容组织逻辑上已完成领域划分，因此能够在源头上减少跨领域内容混入的情况。

\begin{figure}[H]
    \centering 
    \includegraphics[width=0.8\textwidth]{JiaDian.pdf}
    \caption{以“家电”为核心关键词构建的搜索结果页面 }
    \label{fig:appliance_entry}
\end{figure}

\begin{figure}[H] 
    \centering 
    \includegraphics[width=0.8\textwidth]{CaiZhuan.pdf}
    \caption{美妆频道对应的内容流页面 }
    \label{fig:beauty_entry}
\end{figure}

在具体采集过程中，分别从上述不同入口页面解析并提取笔记唯一标识（ID），再基于笔记 ID 进一步访问笔记详情页获取完整内容信息。通过这种“入口区分—ID 提取—详情采集”的两阶段方式，实现了家电与美妆两个内容领域样本的独立采集与明确区分。

不同领域入口URL设置示例:
\begin{lstlisting}[language = Python, numbers=left, 
	numberstyle=\normalsize,keywordstyle=\color{blue!70},showstringspaces=false,
	commentstyle=\color{red!50!green!50!blue!50},frame=shadowbox,
	rulesepcolor=\color{red!20!green!20!blue!20},basicstyle=\ttfamily]
# 本代码仅为过程抽象示例，非实际运行代码。  

# 不同领域入口URL设置
ENTRY_URLS = {
    "家电": (
        "https://www.xiaohongshu.com/search_result/?"
        "keyword=家电&source=web_explore_feed&type=51"
    ),
    "美妆": (
        "https://www.xiaohongshu.com/explore?"
        "channel_id=homefeed.cosmetics_v3"
    )
}
def collect_ids_by_field(field_name):
    """
    根据领域名称选择不同入口URL
    """
    url = ENTRY_URLS[field_name]
    print(f"当前采集领域：{field_name}")
    print(f"使用入口：{url}")
    
    response = session.get(url, headers=get_headers(), cookies=cookies)
    soup = BeautifulSoup(response.text, "html.parser")
    
    # 后续统一执行ID提取逻辑
    # ...

# 调用示例
collect_ids_by_field("家电")
collect_ids_by_field("美妆")


\end{lstlisting}

\subsection{辅助数据源说明}
本研究未引入任何辅助数据源，如微博热搜、百度指数或第三方社交媒体监测平台的数据。所有分析均基于所采集的小红书平台笔记数据展开。
通过构建大规模、高维度的原始笔记数据集，直接从平台内部数据中挖掘影响内容传播效果的关键因素，以避免外部数据噪声对分析结果的干扰。

\section{数据采集方法}
研究采用基于 Python 的自动化网页数据采集方法，对平台公开页面进行结构化解析。
在采集过程中，通过合理控制请求频率与访问行为，尽量模拟真实用户浏览过程，以保证数据采集的稳定性与连续性。-抓取”的流水线作业，提升了系统的健壮性和可维护性。

\subsection{采集工具与框架}
数据采集主要基于 Python 语言实现，使用 requests 库完成网页请求，利用 BeautifulSoup 对 HTML 页面进行解析与信息提取。
通过维护会话状态、设置随机 User-Agent 及请求延时，有效降低高频访问带来的风险，提高采集过程的稳定性。

整体爬虫程序采用模块化设计，分为笔记 ID 采集模块与笔记详情信息采集模块，便于后续维护与扩展。\begin{itemize}
	\item 第一阶段：ID收集爬虫：该模块负责从平台信息流中持续发现新的笔记ID。它模拟访问指定的频道发现页（如美妆、家电频道），解析页面源码，提取笔记详情页链接并剥离出唯一的笔记ID。
	\item 第二阶段：详情页采集爬虫：该模块读取第一阶段生成的ID列表，批量构建详情页URL并发起请求。通过解析详情页的<meta>标签和特定DOM元素，精确提取2.1.1节中定义的各类信息。
\end{itemize}
小红书笔记ID采集流程示例：
\begin{lstlisting}[language = Python, numbers=left, 
	numberstyle=\normalsize,keywordstyle=\color{blue!70},showstringspaces=false,
	commentstyle=\color{red!50!green!50!blue!50},frame=shadowbox,
	rulesepcolor=\color{red!20!green!20!blue!20},basicstyle=\ttfamily]
import time
import random
import requests
from bs4 import BeautifulSoup

def get_note_ids_from_page(url):
    headers = {
        "User-Agent": "Mozilla/5.0"
    }
    response = requests.get(url, headers=headers)
    soup = BeautifulSoup(response.text, "html.parser")

    note_ids = []
    for link in soup.find_all("a", href=True):
        if "/explore/" in link["href"]:
            note_id = link["href"].split("/explore/")[-1]
            note_ids.append(note_id)

    time.sleep(random.uniform(2, 5))  # 随机延时
    return note_ids
\end{lstlisting}

在获取笔记唯一标识后，进一步访问笔记详情页面，对笔记的文本内容、互动数据及发布时间等关键信息进行解析。其基本实现流程如代码示例所示：
\begin{lstlisting}[language = Python, numbers=left, 
	numberstyle=\normalsize,keywordstyle=\color{blue!70},showstringspaces=false,
	commentstyle=\color{red!50!green!50!blue!50},frame=shadowbox,
	rulesepcolor=\color{red!20!green!20!blue!20},basicstyle=\ttfamily]
import requests
from bs4 import BeautifulSoup
import time
import random

def get_note_detail(note_id):
    url = f"https://www.xiaohongshu.com/explore/{note_id}"
    headers = {
        "User-Agent": "Mozilla/5.0"
    }

    response = requests.get(url, headers=headers)
    soup = BeautifulSoup(response.text, "html.parser")

    # 提取笔记标题
    try:
        title = soup.select_one('meta[name="og:title"]')['content']
    except:
        title = ""

    # 提取正文描述
    try:
        description = soup.select_one('meta[name="description"]')['content']
    except:
        description = ""

    # 提取互动数据
    try:
        like_count = soup.select_one('meta[name="og:xhs:note_like"]')['content']
        comment_count = soup.select_one('meta[name="og:xhs:note_comment"]')['content']
        collect_count = soup.select_one('meta[name="og:xhs:note_collect"]')['content']
    except:
        like_count, comment_count, collect_count = 0, 0, 0

    # 提取发布时间
    try:
        publish_time = soup.select_one('span.date').text.strip()
    except:
        publish_time = ""

    time.sleep(random.uniform(2, 4))  # 控制访问频率

    return {
        "note_id": note_id,
        "title": title,
        "description": description,
        "like": like_count,
        "comment": comment_count,
        "collect": collect_count,
        "publish_time": publish_time
    }
\end{lstlisting}

\subsection{采集流程设计}
本研究的数据采集流程主要包括以下几个步骤：
\begin{enumerate}
	\item \textbf{确定采集范围：} 明确研究领域为家电与美妆两个内容类别，数据采集时间范围限定为 2023 年 1 月至 2024 年 8 月。
	\item \textbf{笔记 ID 获取：} 分别通过搜索页面与频道页面获取笔记列表，从页面中解析并提取笔记唯一标识（ID），并对已采集 ID 进行去重处理。
	\item \textbf{笔记详情采集：} 基于已获取的笔记 ID，进一步访问笔记详情页面，采集完整的文本内容、互动数据、多媒体信息及博主相关信息。
	\item \textbf{数据存储：} 将结构化数据统一存储为 CSV 文件，便于后续统计分析；同时将笔记正文、标签词、视频链接等文本信息以及图片资源分别保存，用于辅助特征提取和描述性分析。
\end{enumerate}
数据采集过程为连续多日运行完成，通过分轮采集与断点续爬机制，保证在长时间运行条件下的数据完整性与稳定性。

\subsection{反爬策略与伦理考量}
为降低对平台服务器的访问压力并遵守网络爬虫伦理，本研究在数据采集过程中采取了以下措施：
\begin{itemize}
	\item 设置随机请求延时，避免高频访问；
	\item 控制单轮采集规模，防止异常请求行为；
	\item 仅采集平台公开展示的数据内容；
	\item 所有数据仅用于学术研究，分析结果以聚合统计形式呈现，不涉及任何个人隐私信息的公开。
\end{itemize}本研究所采集的全部数据均为小红书平台上的公开信息。在数据使用上，我们郑重承诺：所有数据仅用于本项学术研究的统计分析，绝不会用于任何商业目的，也不会以任何形式公开披露能识别到具体用户或个人身份的原始信息。论文中的所有结论均基于聚合后的统计结果呈现。

\section{数据描述与特征工程}

原始采集的小红书笔记数据以文本和半结构化字段为主，包含大量非标准化表达形式，无法直接用于统计分析与建模。因此，本研究在数据清洗的基础上，围绕时间信息、互动行为、发布地点与标签等核心要素，构建了一套结构化特征体系，将原始字段转化为可量化、可编码的分析变量。

\subsection{时间特征构建}

小红书平台中的发布时间以多种自然语言形式呈现，如“3天前”“昨天 20:15”“06-20”“编辑于 今天”等。针对该特点，本文对发布时间字段进行统一解析与推断，主要步骤包括：

\begin{itemize}
\item 识别相对时间表达（如“x天前”“昨天”“今天”），并基于数据采集日期反推出实际发布日期；
\item 对仅包含“月-日”或“x月x日”的时间表达，结合采集时间补全年份信息；
\item 对已为标准日期格式（如 YYYY-MM-DD）的记录直接保留。
\end{itemize}

经过处理后，所有笔记的发布日期被统一转换为标准日期格式（YYYY-MM-DD），并在此基础上进一步派生出以下时间特征：

\begin{itemize}
\item \textbf{发布日期：} 推断后的笔记实际发布时间；
\item \textbf{星期特征：} 根据发布日期计算星期几，并编码为数值型变量，用于刻画发布时点的周期性特征。
\end{itemize}

\subsection{互动行为特征}

互动行为是衡量笔记热度的核心信息来源，包括点赞数、评论数、收藏数以及发布者粉丝数。然而，这些字段在原始数据中通常以非标准数值形式出现，如“1.2万”“3万+”“10+”等。

为保证数据的可计算性，本文对上述字段进行统一数值化处理：

\begin{itemize}
\item 将带有“万”“万+”等单位的文本数值转换为对应的整数数量级；
\item 对带“+”的数值，在保留数量级信息的基础上进行近似修正；
\item 最终生成相对准确的点赞数、评论数、收藏数及粉丝数数值特征。
\end{itemize}

该处理过程在一定程度上缓解了平台展示规则带来的信息模糊问题，使互动数据能够直接参与后续统计分析与建模。

\subsection{发布地点特征}

发布地点信息来源于笔记发布时间字段中附带的文本描述。本文通过正则匹配与字符串解析方法，从相关字段中抽取用户填写的发布地点信息，并构建地点特征。

在此基础上，采用字典映射方式对地点文本进行编码处理：

\begin{itemize}
\item 对已识别的发布地点建立全局唯一编号映射表；
\item 对同时包含多个地点的情况，保留多值编码结果；
\item 对未填写或无法识别发布地点的笔记，设置统一的缺失标记。
\end{itemize}

此外，为便于模型区分地点信息的完整性，额外构建“是否填写发布地点”的辅助二元特征。

\subsection{标签特征}

笔记中的话题标签反映了内容的主题属性。本文对笔记中包含的标签词进行统计与整理，并采用与发布地点相同的字典化编码策略：

\begin{itemize}
\item 统计全体数据中出现过的标签词集合；
\item 构建标签词到数值 ID 的映射关系；
\item 对缺失标签或无标签笔记设置统一缺失编码；
\item 构建“是否填写标签”的辅助特征。
\end{itemize}

通过上述处理，原始文本标签被转化为可用于建模的离散类别特征。

\section{数据预处理与清洗}

原始采集的数据在结构和内容层面均存在噪声与不一致问题。为保证数据质量，本文在特征工程之前实施了系统的数据预处理与清洗流程。

\subsection{异常与冗余数据清理}

在数据整理过程中，主要执行以下清洗操作：

\begin{itemize}
\item 删除在合并或追加过程中产生的重复表头行；
\item 剔除关键字段（如粉丝数）中包含明显异常字符串的记录；
\item 对无法解析发布时间或缺失核心标识信息的样本予以剔除。
\end{itemize}

\subsection{缺失值处理策略}

考虑到社交平台数据中缺失现象普遍存在，本文未采用简单均值或众数填补，而是通过显式标注缺失状态的方式进行处理：

\begin{itemize}
\item 对发布地点、标签等非关键特征，使用特定编码值表示缺失；
\item 同时构建“是否填写该信息”的辅助特征，保留缺失本身所蕴含的信息；
\item 对关键字段（如笔记 ID、互动数据）严重缺失的样本直接剔除。
\end{itemize}

\subsection{特征编码与一致性控制}

为保证多文件、多批次数据处理结果的一致性，代码采用增量式字典编码策略：

\begin{itemize}
\item 地点与标签编码字典在处理过程中持续更新并全局共享；
\item 相同文本特征在不同数据文件中始终对应相同的数值编码；
\item 对未匹配到的异常类别设置统一的异常标识值。
\end{itemize}

该策略有效避免了因分批处理导致的特征含义不一致问题，为后续建模提供了稳定的数据基础。

\begin{figure}[H]
    \centering
    % 左图：用\text{}包裹特殊符号，转义下划线
    \subfigure[{%
        \parbox{0.45\textwidth}{\centering 清洗前的原始数据\\
        \footnotesize 同色框对应：红\text{=}标签词，蓝\text{=}粉丝量，绿\text{=}发布时间，紫\text{=}发布地点，橙\text{=}点赞数、评论数、收藏数}%
    }]{
        \includegraphics[width=0.45\textwidth]{before_cleaning.pdf}
        \label{fig:before_cleaning}
    }
    \quad
    % 右图：转义下划线（\_），用\text{}包裹等号
    \subfigure[{%
        \parbox{0.45\textwidth}{\centering 清洗后的结构化数据\\
        \footnotesize 同色框对应：红\text{=}标签编码，蓝\text{=}相对准确粉丝数，绿\text{=}推断发布日期、推翻星期几和对应星期编码，紫\text{=}发布地点编码\text{=}相对准确点赞数、相对准确评论数、相对准确收藏数}%
    }]{
        \includegraphics[width=0.45\textwidth]{after_cleaning.pdf}
        \label{fig:after_cleaning}
    }
    \caption{小红书笔记数据清洗前后对比}
    \label{fig:data_cleaning_comparison}
\end{figure}


\part{Part III: 小红书笔记多维度数据分析与探索}
\chapterimage{chapter_head_2.pdf} % Chapter heading image

\chapter{笔记内容质量分析}

\section{文本内容深度与可读性分析}
本研究从文本长度、可读性与情感三个维度考察小红书笔记内容质量与互动表现之间的关系。由于实际数据处理中，可读性指标仅能对部分样本稳定返回结果，本文采用“全量基础清洗 + 抽样计算特征 + 回归稳健性检验”的方式开展分析。

\subsection{变量构建与样本处理}
首先，本文将原始 CSV 数据合并后进行基础清洗。剔除点赞、评论、收藏缺失的样本后，原始样本量为323,598条，清洗后保留297,031条。为降低极端账号对结果的干扰，本文进一步构造互动率指标，并删除粉丝数异常高但互动率过低的疑似买粉样本，最终样本量仍保持为297,031条。

在核心变量构建上，本文将互动表现定义为：
\[
\text{engagement}=\text{likes}+\text{favs}+4\times \text{comments}
\]
并对其取对数得到以缓解长尾分布对回归估计的影响。文本长度方面，本文计算了标题长度、正文长度以及全文长度，其中全文长度作为主分析变量进入回归模型。情感变量则由情感预测值构造得到，进一步计算情绪极端程度
\[
\text{emotion\_extreme} = \text{emotion\_strength} - 0.5
\]
\subsection{文本长度与互动关系}
为了考察文本长度是否影响用户互动，本文首先对全文长度与互动水平进行了可视化分析。结合后续回归结果，长度的一次项系数为负、二次项系数为正，说明文本长度与互动之间存在一定的 U 型关系。在当前样本区间内，较短文本整体上更容易获得较高互动，但当长度继续增加到一定程度后，互动水平可能再次回升。
\begin{figure}[H]
    \centering 
    \includegraphics[width=0.8\textwidth]{clve.pdf}
    \caption{文本长度与互动关系 }
    \label{fig:content_length_engagement}
\end{figure}

\subsection{文本情感与互动关系}
在情感维度上，本文进一步分析了情感强度与互动表现之间的关系。回归结果表明，情感变量与互动水平呈现负相关趋势。情感强度和情绪极端程度在主模型与稳健性模型中均表现为显著负相关，说明在本样本中，越偏离中性、越强烈的情绪表达，平均互动反而更低。
\begin{figure}[H]
    \centering 
    \includegraphics[width=0.8\textwidth]{eve.pdf}
    \caption{文本情感与互动关系 }
    \label{fig:emotion_engagement}
\end{figure}

\subsection{互动分布特征}
从整体分布来看，笔记互动量呈现明显的长尾特征。大多数样本集中在较低互动区间，只有少量内容获得了极高互动。这也解释了为何本文对互动变量采用对数变换后再进行建模，以提高回归估计的稳定性并减弱极端值的影响。
\begin{figure}[H]
    \centering 
    \includegraphics[width=0.8\textwidth]{edl.pdf}
    \caption{互动特征分布 }
    \label{fig:engagement_distribution}
\end{figure}

\subsection{非线性趋势识别}
为了进一步识别长度变量与互动之间的非线性结构，本文对抽样样本进行了 LOWESS 平滑拟合。结果显示，文本长度与互动之间存在较明显的弯曲趋势，这说明简单线性模型不足以完整描述二者关系。因此，在后续 OLS 回归中，本文加入了长度平方项，以刻画可能存在的非线性效应。
\begin{figure}[H]
    \centering 
    \includegraphics[width=0.8\textwidth]{lowess.pdf}
    \caption{非线性趋势识别 }
    \label{fig:length_lowess}
\end{figure}

\subsection{回归结果与稳健性检验}
在主回归模型中，本文以 
$engagement\_log$ 为因变量，以文本长度、长度平方项、情感强度、情绪极端程度和可读性为自变量进行 OLS 估计。结果显示，模型整体显著，但解释力有限，$R^2$约为0.050。这说明文本特征对互动具有统计意义上的解释力，但笔记热度仍受到账号属性、发布时间、封面质量、主题类别等其他因素影响。

具体而言，文本长度一次项系数显著为负，而二次项系数显著为正，说明长度效应具有非线性特征；情感强度与情绪极端程度的系数均显著为负，表明在本样本中，情绪越强或越极端，平均互动反而越低；可读性变量在主模型中边际显著，但在稳健性模型中不再显著，说明其影响较弱，且对样本极端值较为敏感。

进一步地，本文剔除互动前1%的极端值后重新估计模型，发现核心变量的方向和显著性基本保持一致，说明上述结论具有较好的稳健性。换言之，文本长度与情感特征对互动的影响并非由少数极端爆款样本驱动，而是具有较稳定的总体规律。
\subsection{分组分析}
为更直观地观察长度效应，本文进一步将文本长度按分位数分组，比较不同长度区间的平均互动水平。结果表明，文本长度与平均互动之间存在明显的分组差异，整体趋势与回归结果一致，即长度并非简单线性越长越好，而是存在相对更优的区间。
\begin{figure}[H]
    \centering 
    \includegraphics[width=0.8\textwidth]{ol.pdf}
    \caption{分组分析 }
    \label{fig:length_group_analysis}
\end{figure}

\subsection{情感分组结果}
最后，本文对情感强度与情绪极端程度进行了分组比较。结果表明，情感强度较低、情绪更接近中性的样本平均互动更高；而情绪极端程度越高，平均互动水平总体越低。这一结果与回归分析保持一致，进一步说明在小红书笔记场景中，过强或过于极端的情绪表达并不一定带来更高互动。

综合来看，文本内容特征对笔记热度具有显著影响，其中长度呈现非线性关系，情感强度与情绪极端程度总体上与互动负相关，可读性影响相对较弱且不够稳健。该结果说明，小红书用户对内容的响应不仅取决于文本长度本身，也与内容表达方式、情绪呈现方式及其是否符合平台阅读习惯密切相关。

\section{主题相关性与信息价值分析}

在社交媒体中，除了基础的情感与篇幅特征外，内容传递的实际信息量以及“标题-正文”间的一致性，也是决定用户深度互动的关键因素。本节基于 TF-IDF 算法提取文本核心关键词，从“信息密度”与“标题匹配度”两个维度，深入探讨小红书笔记的信息价值分布对互动表现的影响。

\subsection{变量构建与共线性修正}
为量化笔记的信息分布情况，本文首先利用 TF-IDF 算法提取单篇笔记中的高权重核心词汇。“信息密度（info\_density）”定义为正文高权重词汇数占局部总有效词数的比重；“标题匹配度（title\_match）”则定义为标题关键词与正文关键词的交集词数占标题总词数的比例，主要用于衡量笔记是直接“白描”还是存在一定的内容悬念差。

在将上述变量引入回归模型前，考虑到前期引入的文本长度平方项（$content\_len^2$）会与一次项产生较强的自然数学共线性（未标准化模型的条件数曾高达 $2.04 \times 10^6$），本文利用 Z-score 对所有连续型自变量先行进行标准化处理，再由标准化后的长度构造平方项。处理后，各变量的方差膨胀因子（VIF）介于 1.11 至 7.74 之间，模型总体条件数（Cond. No.）降至 8.52。该结果说明数值尺度导致的病态程度明显缓解，但个别变量仍存在中等程度的共线性，解释单个系数时仍需保持谨慎。

\subsection{信息价值回归结果与检验}
本文将标准化后的各特征变量纳入 OLS 回归模型，因变量为对数互动量 $engagement\_log$，而非除以粉丝数后的互动率。回归估计结果如表 \ref{tab:ols_results_info} 所示。整体模型高度显著（F统计量 = 259.7，Prob(F) $< 0.001$）。

\begin{table}[H]
    \centering
    \caption{信息价值与互动表现的多元线性回归结果}
    \label{tab:ols_results_info}
    \begin{tabular}{lcccc}
        \toprule
        \textbf{变量} & \textbf{标准化系数 ($\beta$)} & \textbf{标准误(Std.Err)} & \textbf{t 值} & \textbf{P>|t|} \\
        \midrule
        截距项 (Intercept)     & 5.4472  & 0.021 & 255.55 & $<$0.001 *** \\
        内容长度 ($content\_len$)      & -0.4994 & 0.043 & -11.71 & $<$0.001 *** \\
        内容长度平方 ($content\_len^2$)& 0.0711  & 0.015 & 4.80   & $<$0.001 *** \\
        情感强度 ($emotion\_strength$) & -0.1100 & 0.018 & -5.97  & $<$0.001 *** \\
        极端情感 ($emotion\_extreme$)   & -0.0969 & 0.019 & -5.18  & $<$0.001 *** \\
        信息密度 ($info\_density$)      & 0.1055  & 0.025 & 4.20   & $<$0.001 *** \\
        标题匹配度 ($title\_match$)     & -0.1822 & 0.016 & -11.23 & $<$0.001 *** \\
        \midrule
        \multicolumn{5}{l}{样本量 = 20,000 \quad $R^2$ = 0.072 \quad 总体条件数(Cond. No.) = 8.52} \\
        \bottomrule
    \end{tabular}
    \vspace{1ex}
    \\
    \small{\textit{注：*** 代表 1\% 的显著性水平。}}
\end{table}

由于引入了标准化处理，相较于前述的基础回归模型，本模型中的各项系数代表了不同特征变量在同一量纲下对互动率的相对作用强度。前述关于长度“U型折返”以及极端情绪负面作用的结论在加入新特征后依旧十分稳健，并且本部分揭示了信息维度两个额外的核心机制。

\subsection{信息密度的“干货溢价”效应}
回归结果表明，信息密度与互动表现呈显著正相关（$\beta=0.1055, p<0.001$）。在控制模型中其他文本变量后，高权重词汇占比较高的笔记通常具有更高的平均互动量。这一结果与“信息更集中、更便于检索的内容可能更符合用户决策需求”的解释一致。不过，本文的信息密度来自TF--IDF阈值构造，且模型尚未控制全部账号与主题差异，因此该结果应理解为相关性证据，而非信息密度对互动的确定性因果效应。

\subsection{预期偏差与“悬念留白”机制}
模型同时揭示了一个值得关注的反直觉现象：标题匹配度与互动表现呈显著负相关（$\beta=-0.1822, p<0.001$），且其标准化系数绝对值在模型中仅次于文本长度。这意味着，在当前样本和词汇重合度定义下，标题关键词被正文直接重复得越多，笔记的平均互动量反而越低。

这一现象可能与信息缺口或标题的悬念表达有关：过分直白的重复可能降低读者继续探索的意愿，而适度的信息差可能保留阅读动机。不过，当前指标衡量的是分词后的词汇重合，而不是严格的语义一致性；标题较短、正文缺失或不同主题的语言习惯也可能影响该指标。因此，“悬念留白”只能作为一种可能解释，尚不能由本模型直接验证。

\subsection{信息分布的可视化特征}
为进一步佐证回归模型揭示的两大规律，本文对信息密度、标题匹配度与互动强度的关系进行了二维六边形密度成像拟合（Hexbin Plot）。如图 \ref{fig:info_hexbin} 所示，左侧图中高密度分布带向右侧高互动区的蔓延直观体现了“干货溢价”的底盘支撑效能；而右侧图则清晰地反映了标题匹配度过高时，样本互动数据出现的集中性坍缩，为本节总结的“悬念留白”机制提供了几何维度的实质性支撑。

\begin{figure}[H]
    \centering 
    \includegraphics[width=0.8\textwidth]{hexbin_density.pdf}
    \caption{信息密度与标题匹配度对互动的二维密度影响分析}
    \label{fig:info_hexbin}
\end{figure}

\chapter{笔记形式与视觉呈现分析}

\section{多媒体内容分析}
小红书是一个强视觉化的平台，图片和视频的质量直接影响用户的点击和停留意愿。本研究将对笔记的多媒体内容进行量化分析。

首先，我们将分析图片的基本特征。对于包含图片的笔记，我们将计算其图片数量，并研究图片数量与互动量的关系。通常，一定数量的高质量图片能更好地展示内容，但过多图片可能导致加载缓慢，反而影响体验。我们还将通过图像处理技术提取封面图或首图的关键视觉特征，如：
\begin{itemize}
	\item \textbf{色彩特征：} 主色调、色彩丰富度（饱和度、对比度）。明亮、色彩和谐的图片可能更具吸引力。
	\item \textbf{构图特征：} 是否包含人脸、人脸占比、是否遵循三分法等构图原则。
	\item \textbf{清晰度与美学质量：} 使用图像清晰度算法（如拉普拉斯方差）和预训练的深度学习美学评估模型对图片质量进行评分。
\end{itemize}
我们将通过相关性分析和分组实验（如对比高热度笔记和低热度笔记的图片特征均值）来验证哪些视觉特征与笔记热度显著相关。

其次，对于视频笔记，我们将分析视频时长、封面帧、是否有字幕、背景音乐等特征。视频时长存在一个“甜点区间”，过短可能表达不清，过长则完播率低。我们将通过数据分析寻找不同内容类别下最优的视频时长范围。

\section{排版与格式分析}
在小红书这类强视觉驱动的社交媒体平台上，笔记的排版与格式直接影响用户的阅读体验和信息获取效率。 良好的排版不仅能够降低用户的认知负荷，还能通过视觉引导突出关键信息，从而延长停留时间并激发互动行为。尽管排版特征往往被视为“微内容 ”，但其对传播效果的潜在贡献不可忽视 。本研究旨在系统量化排版与格式特征对笔记热度的影响，为内容创作者提供可操作的优化建议。
\begin{itemize}
	\item \textbf{文本分段：} 正文是否合理分段（通过换行符或分隔符数量衡量），长段落是否使用小标题或序号。
	\item \textbf{重点突出：} 检测正文中是否包含加粗标签（<b>、<strong>） 、特殊符号或表情符号（Unicode Emoji）来强调关键信息。
	\item \textbf{清单使用：} 检测正文是否包含列表结构，或平台特有的列表格式开头的行，是否使用“1. 2. 3.”或“\textbullet{}”等列表形式，使内容条理清晰。特征通过Python的正则表达式库re和emoji库进行提取。
\end{itemize}
这些格式特征虽然看似细微，但可能通过影响阅读流畅度和信息获取效率，间接影响用户的互动意愿。我们将通过构建二元特征（如“是否使用表情符号”）并比较其在不同热度组中的分布，来评估其影响。

\chapterimage{chapter_head_2.pdf}
\chapter{发布者影响力与社交网络分析}

\section{发布者基础影响力分析}
发布者自身的社交资产是笔记获得初始流量的基石。粉丝数量决定了笔记发布时的基础曝光池大小，而历史表现和认证状态则反映了发布者的内容生产能力和用户信任度。本部分旨在量化发布者基础特征对笔记热度的贡献，并识别高粉丝粘性的“优质创作者”特征。   
\begin{figure}[htbp]
    \centering                % 图片居中
    \includegraphics[width=0.8\textwidth]{16.png}  % 宽度设为页面宽度的80%
    \caption{不同粉丝量级笔记的平均点赞数}  % 图片标题
    \label{fig:followers_likes}       % 标签，供引用
\end{figure}
\begin{itemize}
	\item \textbf{粉丝数量（取对数）：} 从用户信息页采集粉丝数，由于分布呈长尾，对其进行自然对数变换（log(粉丝数+1)）以缓解异方差性。计算粉丝数对数与笔记点赞数对数的皮尔逊相关系数，并通过散点图和平滑拟合（LOESS）观察线性关系强度。
	\item \textbf{历史表现：}计算每个发布者在采集窗口前30天内发布的笔记平均点赞数、评论数、收藏数，作为“近期内容质量 ”的代理变量 。若发布者为新用户，则用整体样本均值填充。定义互动率 = 点赞数 / (粉丝数+1)，计算每个笔记的互动率，并分析其分布 。将发布者按互动率分组， 比较不同组别笔记的平均互动量， 以识别“ 高质量 ”发布者（粉丝数不高但互动率高）。
	\item \textbf{认证状态：} 比较认证用户（企业号、专业号、达人）与非认证用户在平均互动量上的差异。认证通常意味着更高的可信度和专业性，可能带来额外的流量权重。 以对数化后的点赞数为因变量，将粉丝数对数、历史平均表现对数、认证状态等作为自变量，构建多元线性回归模型 。 同时，加入粉丝数与内容特征的交互项，探究不同粉丝量级下内容要素的异质性作用。
\end{itemize}

\section{发布者内容垂直度与粉丝画像（模拟）分析}
由于难以直接获取粉丝画像，我们可以通过分析发布者历史笔记的主题分布来推断其内容垂直度。一个高度垂直的发布者（如只发美妆内容）可能拥有更精准、忠诚的粉丝群体，其笔记的互动率和转化率可能更高。我们将计算发布者历史笔记主题分布的熵值，熵值越低代表垂直度越高，并分析垂直度与平均互动率的关系。

\chapterimage{chapter_head_2.pdf}
\chapter{发布时机与上下文分析}

在社交媒体生态中，内容的传播效能不仅取决于文本与视觉质量，还受到发布时境（Temporal Context）与地域文化（Spatial Context）的深刻影响。本章在利用规则引擎剔除了 15,972 条“高粉丝低互动”的疑似买粉样本后，基于 29.6 万条有效数据，从星期节律、中国特色法定节假日以及地域数字经济三个维度，深度剖析小红书笔记传播的时空动力学机制。

\section{时间节律：基于周周期的注意力波动分析}
受限于现代社会工作与学习的周期性规律，用户的线上注意力呈现出明显的周内波动机制。利用单因素方差分析（ANOVA），检验结果表明“星期几效应对数互动量”的主效应极显著（F = 946.56, $p < 0.001$，$\eta^2 = 0.018$）。

数据呈现出两个值得关注的极值节点：第一，\textbf{周一是内容互动的最高峰（平均 log 互动量达 5.95）}。这一反常识现象可用组织行为学中的“心理补偿机制（Psychological Compensation）”来解释——在结束周末回归工作或学业的首日，用户倾向于利用流式媒体进行碎片化的“赛博摸鱼”以缓解职段压力。其二，\textbf{周五是互动的绝对低谷（平均 log 互动量仅为 5.01）}。Tukey HSD 事后检验进一步表明，周五与其他任何工作日的均值差异（Mean Diff 均在 $-0.5$ 以上）均极其显著（$p_{adj} < 0.01$）。这揭示了媒介使用的替代效应（Substitution Effect）：在周末临近前，受众注意力从“线上内容消费”发生了向“线下实体社交”的物理转移。

\begin{figure}[H]
    \centering 
    \includegraphics[width=\textwidth]{weekday_effect.pdf}
    \caption{小红书笔记互动的周周期波动及分指标分布}
    \label{fig:weekday_effect}
\end{figure}

\section{制度性节点冲击：节假日与调休制度的异象}
除常规节律外，国家法定节假日引起的社会作息变动对平台流量池具有强烈的宏观冲击。为了精准刻画这一效应，本文将 2023-2024 年的数据严格划分为：普通工作日、普通周末、节假日以及中国特有的“调休工作日”。

四组样本的 ANOVA 检验显示出极显著的结构差异（F = 305.53, $p < 10^{-190}$）。其中最引人注目的实证发现是：\textbf{调休工作日的互动峰值反常地超越了所有日期类型，平均总互动跃升至 4744（对数互动 6.28），甚至高于真正的法定节假日（4202）。}

一方面，节假日（Welch t = 18.75, $p < 0.001$ 显著高于非节假日）作为大块闲暇时间，确实是深度长决策内容（如旅游攻略、家电选购）消费的集中区。但另一方面，调休工作日的极高数据表征了“强迫性劳动时间”与“注意力涣散”的错位：处于心理疲劳与调休补偿心态下的用户，其移动端浏览频率剧增。在该节点发布的内容，极易借由宏观“摸鱼流量池”脱颖而出。

\begin{figure}[H]
    \centering 
    \includegraphics[width=0.8\textwidth]{holiday_effect.pdf}
    \caption{不同日期类型（含调休特征）对互动表现的异质性冲击}
    \label{fig:holiday_effect}
\end{figure}

\section{空间格局：地域数字文化的非对称扩张}
从数字地理学的视角切入，社交媒体的数据特征往往能够折射出区域经济水平与消费文化的图景。在剔除无效地理标记后，本文梳理了样本分布排名前列的省份（市），并观察到了“规模优势”与“传播效率高地”并非完全重合的现象。

从总量集聚来看，\textbf{广东省以 25,624 条记录确立了断档第一的大盘领跑属性}（超第二名几乎一倍）。这印证了珠三角庞大的电商与图文代运营产业链造就的内容生产渗透率。然而，从互动转化效率来看，其平均对数互动（5.23）却处于腰部水平，体现了过度商业化与内卷带来的流量互相稀释困境。

另一方面，\textbf{江浙沪与四川盆地构筑了“高生产质感+高互动黏性”的传播堡垒}。浙江（5.97）、四川（5.88）与上海（5.87）在拥有巨大基数的同时，维持了全局最高梯队的绝对互动均值。这在宏观上对应了此类核心地带用户对其输出的“精致生活方式（Lifestyle）”以及活跃“松弛社交文化（Chill Culture）”的向外辐射与强认同。

\begin{figure}[H]
    \centering 
    \includegraphics[width=\textwidth]{province_analysis.pdf}
    \caption{小红书内容互动效能的地域热度解构}
    \label{fig:province_analysis}
\end{figure}

\chapterimage{chapter_head_2.pdf}
\chapter{标签策略与话题热度分析}

\section{标签使用策略分析}
话题标签是小红书内容分类、搜索曝光和推荐分发的重要工具。合理使用标签能够将笔记精准推送给目标受众，同时借助热门标签获取泛流量。本部分旨在系统研究标签数量、类型组合及标签热度对笔记传播效果的影响，为创作者提供科学的标签策略指导。
\begin{figure}[htbp]          % h=here, t=top, b=bottom, p=page
    \centering                % 图片居中
    \includegraphics[width=0.8\textwidth]{17.png}  % 宽度设为页面宽度的80%
    \caption{不同标签数量笔记的平均点赞数}  % 图片标题
    \label{fig:tag_count_likes}       % 标签，供引用
\end{figure}
\begin{itemize}
	\item \textbf{标签数量：} 统计笔记使用的话题标签数量，分析其与热度的关系。过多标签可能被视为spam，过少则可能降低曝光机会。统计每篇笔记包含的标签个数，构建连续特征。绘制标签数量的分布直方图，并统计不同标签数量区间下笔记的平均互动量，寻找最优标签数量区间。
	\item \textbf{标签类型组合：}通过预定义的标签词库（基于数据中高频标签的聚类）将标签划分为以下类型：通用热门标签、垂直细分标签、地点标签、品牌标签等等。对每篇笔记，统计各类标签的数量。分析高热度笔记的标签组合策略，例如是否同时使用了“大流量泛标签”和“精准垂类标签”。使用关联规则挖掘或决策树模型，分析高热度笔记中标签类型的共现模式，识别典型的“黄金标签组合 ”。
	\item \textbf{标签热度：} 计算每个标签在数据采集期间出现的总频次，作为其“基础热度”。分析笔记所带标签的平均热度或最高热度是否与笔记自身热度相关。将笔记按标签平均热度分位数分组， 比较各组互动量的差异，并通过回归模型评估在控制粉丝数等因素后标签热度的边际贡献。

\end{itemize}

\section{话题借势与热点响应分析}
社交媒体内容传播具有显著的时效性，紧跟热点事件是获得爆发式流量的常见策略。然而，盲目“蹭热点”可能导致内容与账号定位不符，甚至引发用户反感。本部分旨在量化热点响应对笔记热度的实际增益，并探索最佳介入时机，为内容营销的时效性决策提供数据支撑。
\begin{itemize}
	\item 效果对比：采用倾向得分匹配（PSM）或逆概率加权（IPW）等方法，为每篇“蹭热点”笔记匹配一组特征相似的非热点笔记，比较两组在发布后72小时内的平均互动量差异，从而得到热点响应的净效应。
	\item 时机效应分析：将热点笔记按介入阶段分组，使用方差分析比较各组互动量的均值，并通过回归模型控制其他因素，估计不同阶段的边际收益。同时，引入热点类型（节日型、事件型、娱乐型）与阶段的交互项，探究异质性。
	\item 生命周期建模：对每个热点事件，追踪其相关笔记的互动量随时间衰减曲线，使用生存分析或时序模型拟合热点的影响力衰减速度，为预测热点剩余价值提供依据。
\end{itemize}
对于高度相关的特征对（如相关系数绝对值大于0.8），在后续建模中可能需要考虑只保留其中一个或进行降维处理。

\chapterimage{chapter_head_2.pdf}
\chapter{多维度相关性分析}

\section{特征间相关性分析}
为避免原始文件、筛选文件及中间结果重复计入，本章仅读取文件名中包含“笔记详细信息”的数据，并优先按照笔记链接去重；链接缺失时，采用编号、标题和博主名称的组合去重。完成互动指标清洗后，共得到182,103条可用笔记，其中168,323条具有有效粉丝数。由于现有数据不包含可靠的图片数量、认证状态和发布小时，本章不对这些变量进行推断，而使用实际可获得的粉丝数、文本长度、标签数、标点与表情符号、情感得分、标题正文重合度、星期和地点完整性等特征。

考虑到互动量和粉丝数均呈明显长尾分布，本文分别构造
\[
 y=\log(1+\text{engagement}),\qquad
 x_{fans}=\log(1+\text{fans}),
\]
并同时计算Pearson相关系数与Spearman秩相关系数。前者用于识别线性关系，后者对极端值更稳健，可用于判断变量之间是否存在稳定的单调关系。图\ref{fig:chapter8_pearson}给出了特征间Pearson相关系数矩阵。

\begin{figure}[H]
    \centering
    \includegraphics[width=0.92\textwidth]{chapter8_pearson_heatmap.png}
    \caption{数值特征间的Pearson相关系数矩阵}
    \label{fig:chapter8_pearson}
\end{figure}

结果显示，正文长度与全文长度的相关系数达到$r=0.988$，是唯一满足$|r|\geq0.8$的特征对。这一高度相关主要源于全文长度由标题和正文共同构成，属于变量定义带来的机械相关。在后续多变量模型中不宜同时保留正文长度与全文长度，本文建议保留含义更完整的全文长度。其余变量之间未出现绝对相关系数超过0.8的组合，说明除长度变量外，当前特征集合不存在严重的两两线性共线性。需要强调的是，两两相关较低并不能完全排除多变量共线性，后续回归仍应结合VIF进行诊断。

\section{特征与目标变量相关性分析}
本文首先以对数互动量为连续型目标，分别计算各特征的Pearson和Spearman相关系数，并采用Benjamini--Hochberg方法对多重检验的$p$值进行FDR校正。图\ref{fig:chapter8_target}按照Spearman相关系数绝对值展示变量排序。

\begin{figure}[H]
    \centering
    \includegraphics[width=0.82\textwidth]{chapter8_target_ranking.png}
    \caption{数值特征与对数互动量的Spearman相关系数排序}
    \label{fig:chapter8_target}
\end{figure}

粉丝数对数与对数互动量的相关性最强，Pearson相关系数为0.489，Spearman相关系数为0.469，均在FDR校正后显著。这表明发布者既有受众规模仍是笔记互动表现最重要的单一关联因素，但相关系数未达到高度相关水平，笔记热度并不能仅由粉丝基数解释。正文长度和全文长度与互动量呈弱负相关，Spearman系数分别为$-0.116$和$-0.115$；是否具有有效发布地点的系数为$-0.111$。情绪极端程度、标题正文重合度和正向情感概率的Spearman系数分别为$-0.070$、$-0.068$和$-0.063$，方向与第3章的回归结果基本一致，但效应强度较小。标签数量、表情符号、星期、感叹号和问号等变量的绝对相关系数均小于0.05，统计显著性主要受大样本量影响，不宜解释为具有较强的实际作用。

为进一步比较高热度内容，本文将总互动量位于样本前10\%的笔记定义为“爆款”，并对连续特征实施点二列相关、Welch独立样本$t$检验和Mann--Whitney U检验，同时报告Cohen's $d$。粉丝数对数的组间差异最明显：普通组均值为7.613，爆款组均值为9.589，$d=0.696$，属于中等偏大的效应。相比之下，爆款组正文长度和全文长度更短，对应效应量分别为$-0.213$和$-0.211$；标签数量、情绪极端程度与表情数量的效应量绝对值约为0.12至0.14，其余变量的效应量更小。问号数量和星期编码在FDR校正后均未表现出稳定的爆款组差异。

星期编码本身不适合作为严格连续变量解释，因此本文另行进行分组描述。如图\ref{fig:chapter8_weekday}所示，周六和周日的爆款比例分别为10.91\%和10.86\%，周一为10.42\%，周五最低，为8.54\%。该结果说明星期差异主要体现为类别间波动，而非随编码增加或减少的线性趋势。

\begin{figure}[H]
    \centering
    \includegraphics[width=0.9\textwidth]{chapter8_weekday.png}
    \caption{不同星期的互动量分布与爆款比例}
    \label{fig:chapter8_weekday}
\end{figure}

总体而言，粉丝规模是与互动量及爆款状态关联最稳定的特征；文本长度表现出弱负相关；情感、标签及格式变量虽然在大样本下多达到统计显著，但实际效应普遍较弱。因此，后续预测模型不应仅根据$p$值筛选变量，而应综合考虑效应量、非线性结构与变量的理论意义。

\section{非线性关系探索}
线性相关系数可能掩盖阈值、平台期或转折关系。为此，本文对粉丝数对数、全文长度、标签数量、正向情感概率、情绪极端程度和标题正文重合度进行LOWESS局部加权平滑，并将全部有效样本划分为20个分位数组计算组内平均互动。为减少极端值对曲线边界的影响，作图时剔除各自变量上下0.5\%的观测；散点和LOWESS最多随机抽取20,000条记录，而分位数组均值仍使用全部有效数据。

图\ref{fig:chapter8_lowess_main}展示了粉丝规模与文本长度的主要非线性形态。粉丝数对数与互动量总体呈明显正向关系，但曲线并非全程保持相同斜率，说明粉丝增长带来的互动增量存在阶段差异。全文长度曲线则表现出较明显的弯曲：短文本区间互动水平相对较高，随着长度增加先下降，随后在部分较长文本区间趋于平缓或出现局部回升。这与第3.1节在OLS中加入长度平方项后得到的非线性结论相互印证。

\begin{figure}[H]
    \centering
    \begin{minipage}{0.49\textwidth}
        \centering
        \includegraphics[width=\textwidth]{chapter8_lowess_fans.png}
    \end{minipage}
    \begin{minipage}{0.49\textwidth}
        \centering
        \includegraphics[width=\textwidth]{chapter8_lowess_content.png}
    \end{minipage}
    \caption{粉丝规模与全文长度的LOWESS非线性关系}
    \label{fig:chapter8_lowess_main}
\end{figure}

情绪极端程度的LOWESS结果见图\ref{fig:chapter8_lowess_emotion}。总体趋势显示，越偏离中性区间，平均对数互动量越低，但局部区间存在波动。这说明情感效应不宜简单表述为严格线性下降，更适合在后续模型中使用二次项、限制性立方样条或树模型检验。标签数量、正向情感概率和标题正文重合度也呈现不同程度的局部弯曲，但其总体变化幅度明显小于粉丝规模和文本长度。

\begin{figure}[H]
    \centering
    \includegraphics[width=0.78\textwidth]{chapter8_lowess_emotion.png}
    \caption{情绪极端程度与互动量的LOWESS关系}
    \label{fig:chapter8_lowess_emotion}
\end{figure}

需要指出，LOWESS属于探索性可视化方法，曲线形态不能直接证明因果关系。账号粉丝基础、内容类别、采集时间以及平台推荐机制都可能同时影响文本特征和互动结果。因此，本节的非线性发现主要用于为后续模型设定提供依据，而不将曲线转折点直接解释为普遍适用的内容运营最优值。

\chapterimage{chapter_head_2.pdf}
\chapter{特征选择与重要性评估}

\section{基于统计检验的初步筛选}
在构建复杂的机器学习模型之前，通过统计检验对大量原始特征进行初步筛选 , 可以剔除与目标变量无关或相关性极弱的特征， 降低维度、减少噪声，并提高后续建模的效率与可解释性 。本部分将根据任务类型（回归或分类）采用不同的统计检验方法。
\begin{itemize}
	\item \textbf{针对回归任务（预测互动量）：} 数值型特征：对每个数值特征x，计算其与目标变量y（通常取对数化后的点赞数） 的皮尔逊相关系数r及显著性p值。若p < 0.05，则认为该特征与目标存在显著线性相关，予以保留；否则剔除。由于皮尔逊相关对线性关系敏感，对于可能呈非线性关系的特征（如标签数量） ，可同时计算斯皮尔曼秩相关系数作为补充。
分类型特征：对于类别型特征，使用单因素方差分析（ANOVA） 检验不同类别下目标变量均值是否存在显著差异 。若F检验的p < 0.05 ，则保留该特征；否则剔除。对于类别数较多的特征，可先合并低频类别后再进行检验。
	\item \textbf{针对分类任务（预测是否爆款）：} 数值型特征：将样本按爆款标签分为两组 ，首先进行正态性检验（Shapiro- Wilk或Kolmogorov-Smirnov） 。若两组数据均近似正态分布，则使用独立样本t检验比较均值差异；若任一组的分布非正态 ，则使用Mann-Whitney U检验比较中位数差异。保留p < 0.05的特征。
分类型特征：构建特征与爆款标签的列联表 ，进行卡方独立性检验 。若p < 0.05，则拒绝原假设（特征与爆款独立），保留该特征。对于有序分类特征（如星期几），可同时进行趋势卡方检验。
\end{itemize}

\section{基于机器学习模型的特征重要性评估}
统计检验主要捕捉特征与目标之间的线性或单调关系，而机器学习模型能够自动学习复杂的非线性模式及特征间的交互作用。通过多种模型的内部机制，我们可以获得更全面、更稳健的特征重要性排序，为最终构建可解释的预测模型和提出优化策略提供依据。

\begin{itemize}
	\item \textbf{树模型的重要性：} 随机森林（Random Forest）：训练回归或分类随机森林模型，获取基于基尼不纯度减少（或均方误差减少）的特征重要性评分。该评分衡量了每个特征在树节点分裂时减少不纯度的累计贡献。通过交叉验证选择树的数量和最大深度等超参数；梯度提升树：训练XGBoost或LightGBM模型，获取基于分裂增益（gain）的特征重要性，即每个特征在每次分裂时带来的平均损失减少量。该指标比基于分裂次数的权重更能反映特征的预测能力。
	\item \textbf{正则化线性模型：} Lasso回归（L1正则化）：对特征进行标准化后，训练Lasso回归模型。Lasso通过L1惩罚项将部分特征的系数压缩为零，从而实现嵌入式特征选择。保留系数非零的特征，其绝对值大小反映了特征对预测的贡献程度。由于Lasso对线性关系敏感，可与树模型的结果互补。
	\item \textbf{模型不可知解释方法：} SHAP值：基于博弈论中的Shapley值，为每个样本的每个特征分配一个贡献值，表示该特征对预测结果的边际贡献。计算所有样本的SHAP值绝对值平均值，得到全局特征重要性排序。SHAP值不仅能反映重要性，还能通过依赖图展示特征与预测值的关方向及非线性模式。对于XGBoost或随机森林模型，可使用TreeExplainer高效计算SHAP值。
\end{itemize}


将上述三种方法得到的重要性评分进行归一化处理，并采用加权平均或排名平均的方法获得综合重要性得分。设定阈值，筛选出核心特征集。同时，观察不同方法的一致性：若某个特征在多种方法中均排名靠前，则其重要性更为可信。


\part{Part IV: 基于机器学习的笔记热度分类预测模型}
\chapterimage{chapter_head_2.pdf} % Chapter heading image

\chapter{基准分类模型}

\section{预测问题描述}
在小红书内容传播效果研究中，一个核心任务是根据笔记的多维度特征（如内容、形式、发布者、时间、标签等）预测其是否具有高互动潜力。我们将此任务定义为二分类问题：给定一篇笔记的特征向量，模型需要预测其属于“高热度”（如互动量位于前20\%的爆款）还是“普通热度”类别。这比直接预测连续互动量更稳健，能有效处理极端值并服务于实际的二元决策（如是否加大推广）。通过构建和比较多种经典分类模型，我们可以识别出对预测贡献最大的特征，并理解各因素与热度之间的复杂关系，为内容优化提供可解释的洞见。

\section{逻辑回归}
逻辑回归（Logistic Regression）是解决二分类问题最基础的线性模型，以其简单、高效和高可解释性著称。尽管其假设特征与对数几率间存在线性关系，限制了其捕捉复杂非线性模式的能力，但作为基准模型，它能快速提供特征重要性的初步排序（通过系数大小和方向），帮助我们理解各因素对成为爆款的边际贡献。在处理小红书笔记分类时，逻辑回归可首先用于验证特征工程的合理性，并为后续复杂模型提供性能比较的基准。
\begin{itemize}
	\item \textbf{模型原理：} 通过Sigmoid函数将特征的线性组合映射到(0,1)区间，输出样本属于正类（高热度）的概率。决策边界为线性超平面。
	\item \textbf{优势：} 计算速度快，模型参数（系数）直接可解释，能直观看出各特征对结果的正/负向影响及相对强度。
	\item \textbf{局限：} 无法自动捕捉特征间复杂的交互作用和非线性关系，对数据线性可分性要求较高。
\end{itemize}

\section{支持向量机}
支持向量机（Support Vector Machine, SVM）通过寻找最大化分类间隔的最优超平面来实现分类，在处理高维、非线性问题上表现优异。对于小红书笔记数据，其高维特征（如文本向量、多标签）和潜在的非线性关系（如图片数量与热度的倒U型关系）使SVM成为一个强有力的基准模型。
\begin{itemize}
	\item \textbf{核技巧：} 通过引入高斯径向基函数（RBF）等核函数，SVM能将原始特征空间映射到更高维度，从而在变换后的空间中找到一个线性分界超平面，有效处理原始空间中的非线性模式。
	\item \textbf{优势：} 在高维空间中表现良好，泛化能力强，尤其适用于样本量不是特别巨大的情况。通过调整正则化参数C和核函数参数，可以灵活控制模型复杂度。
	\item \textbf{应用：} 在小红书笔记分类中，SVM可以学习到内容质量、发布者影响力、发布时间等特征之间复杂的相互作用，以区分爆款与普通笔记。
\end{itemize}

\section{集成学习模型}
集成学习通过结合多个基学习器的预测结果，获得比单一模型更稳定、更准确的强学习器，非常适合处理小红书笔记分类中特征多样、模式复杂的问题。

\subsection{随机森林}
随机森林（Random Forest）通过自助采样（Bootstrap）构建多棵决策树，并综合所有树的投票结果进行分类。其引入的随机特征选择有效降低了模型的方差，防止过拟合。
\begin{itemize}
	\item \textbf{优势：} 能处理高维数据，对缺失值不敏感，无需复杂的数据预处理。模型可提供特征重要性评分，有助于我们识别影响热度的关键因素（如封面图质量、特定话题标签）。
	\item \textbf{应用：} 在小红书场景下，随机森林可以同时评估文本、视觉、上下文等数百个特征，并给出重要性排序，例如可能发现“发布时段”和“标题情感”的交互作用对热度预测至关重要。
\end{itemize}

\subsection{梯度提升决策树}
梯度提升决策树（Gradient Boosting Decision Trees, GBDT）及其高效实现（如XGBoost, LightGBM）以串行方式迭代训练决策树，每一棵新树都致力于纠正前一棵树的残差。
\begin{itemize}
	\item \textbf{优势：} 预测精度通常高于随机森林，能自动处理特征的非线性关系和交互作用，且通过正则化控制模型复杂度。XGBoost和LightGBM在计算效率和内存使用上进行了大量优化，适合处理大规模数据集。
	\item \textbf{应用：} 对于拥有海量笔记数据的小红书研究，LightGBM能快速训练并高效处理类别型特征（如发布城市、星期几）。其输出的特征重要性（基于分裂增益或覆盖度）能深度揭示各因素对笔记热度的贡献，例如可能精确量化“使用特定热门标签”对互动量的提升幅度。
\end{itemize}

\section{深度学习模型}
深度学习模型能自动从原始或初级特征中学习高级表示，特别适合处理小红书笔记中的非结构化数据（如文本、图像）。

\subsection{多层感知机}
多层感知机（Multilayer Perceptron, MLP）是最基本的深度学习模型，由全连接层和非线性激活函数堆叠而成。它可以作为深度学习方法的入门基准，用于验证经过特征工程的结构化数据是否包含足够预测信息。
\begin{itemize}
	\item \textbf{优势：} 能够学习特征间的复杂非线性组合，相比逻辑回归具有更强的表示能力。
	\item \textbf{应用：} 可将经过预处理和特征工程的所有数值和类别特征（经过编码后）输入MLP，作为与其他更复杂深度学习模型对比的基线。
\end{itemize}

\subsection{卷积神经网络}
卷积神经网络（Convolutional Neural Network, CNN）在图像处理领域取得巨大成功，但其局部感知和参数共享的思想也可应用于一维序列数据。
\begin{itemize}
	\item \textbf{应用：} 对于文本数据，可将词嵌入序列视为一维“图像”，使用一维CNN来提取短语级别的局部语义特征。对于图像特征，可以直接使用预训练的CNN（如ResNet）提取的深度特征向量作为输入的一部分。CNN能有效捕捉如“标题中的特定关键词组合”或“封面图中同时出现人脸和产品”等局部模式对热度的信号。
\end{itemize}

\subsection{循环神经网络}
循环神经网络（Recurrent Neural Network, RNN）及其变体长短期记忆网络（LSTM）专门为序列数据设计，能捕捉时间或顺序上的依赖关系。
\begin{itemize}
	\item \textbf{应用：} 在小红书笔记分类中，LSTM可用于建模文本的序列信息（如正文的语义流），或建模用户历史行为序列（如发布者过去笔记的互动趋势）以捕捉其内容质量的时序演变。然而，对于单篇笔记的静态特征分类，LSTM的优势可能不如CNN或MLP明显，除非考虑时间上下文（如当前热点话题）。
\end{itemize}

\section{实证设置与基准结果}
为将上述模型由方案说明落实为可复现实验，本文读取工作区三个业务数据压缩包中的笔记详情表，合并后共得到323,598条原始记录。按笔记链接（无链接时按编号、标题和博主）去重并剔除互动指标缺失或为负的记录后，保留199,473条有效笔记。互动量定义为“点赞数+收藏数+$4\times$评论数”，达到全样本第80百分位（2,134）的笔记记为高热度，正类比例为20.01\%。建模输入仅使用发布前或内容本身可获得的10项特征，包括粉丝量对数、标题和正文长度、标签数、感叹号数、问号数、Emoji数、情感得分、星期和是否标注发布地；点赞、收藏和评论没有作为输入，以避免目标泄漏。

考虑计算成本，从有效数据中按类别分层抽取60,000条用于比较，再按8:2分层随机划分训练集与测试集。数值缺失以训练集的中位数填补；逻辑回归、线性SVM和MLP同时进行标准化。逻辑回归与SVM使用类别权重平衡，随机森林使用300棵树，梯度提升采用直方图梯度提升分类器，MLP包含64和32个神经元的两个隐藏层并使用早停。表\ref{tab:chapter10-baselines}报告固定测试集结果。鉴于正类仅占约20\%，模型选择以PR-AUC为主，ROC-AUC、F1、精确率和召回率为辅。

\begin{table}[htbp]
\centering
\caption{第10章基准分类模型测试集结果}
\label{tab:chapter10-baselines}
\begin{tabular}{lccccc}
\toprule
模型 & ROC-AUC & PR-AUC & F1 & 精确率 & 召回率 \\
\midrule
逻辑回归 & 0.700 & 0.405 & 0.414 & 0.303 & 0.649 \\
线性SVM & 0.700 & 0.406 & 0.413 & 0.302 & 0.651 \\
随机森林 & 0.725 & 0.424 & 0.373 & 0.479 & 0.305 \\
梯度提升树 & \textbf{0.736} & \textbf{0.440} & 0.263 & \textbf{0.600} & 0.169 \\
多层感知机 & 0.718 & 0.423 & 0.299 & 0.562 & 0.203 \\
\bottomrule
\end{tabular}
\end{table}

梯度提升树的排序能力最佳（PR-AUC=0.440，ROC-AUC=0.736），说明结构化特征与热度之间存在明显的非线性关系；随机森林和MLP次之。逻辑回归与线性SVM的PR-AUC接近，但在类别加权后召回率均约为0.65，更适合“尽量不漏掉潜在爆款”的初筛场景。梯度提升树默认阈值下精确率最高而召回率较低，若用于候选内容筛选，应在验证集上按业务成本降低分类阈值，而不能只依据默认0.5阈值下的F1作结论。

本轮完成了10.1至10.5.1的可运行实证。10.5.2的CNN与10.5.3的RNN/LSTM未给出数值结果：现有压缩包没有与每条笔记稳定对应的原始图片张量或发布者连续历史序列，直接把静态汇总特征重塑为“图像”或伪序列会制造不具业务含义的结构。虽然正文文本可形成词序列，但在没有独立时间外测试和文本模型调参预算的条件下，其结果不能与本节结构化基准作公平比较。因此本研究保留CNN与RNN的模型设计说明，待补齐图片、严格的用户历史序列和时间切分后再实施，避免报告不可复现或误导性的深度模型成绩。

\chapter{针对小红书数据特点的模型改进方向}

在基准模型的基础上，针对小红书笔记数据高维、稀疏、类别不平衡及多模态融合的特点，我们将探索以下改进方向以提升分类性能。

\section{处理类别不平衡}
小红书笔记热度分布高度不均衡，绝大多数笔记互动平平，仅少数成为爆款。直接使用原始数据训练会使模型严重偏向多数类。我们将采用多种策略：
\begin{itemize}
	\item \textbf{数据层重采样：} 使用SMOTE（Synthetic Minority Oversampling Technique）为少数类（高热度笔记）合成新样本，或使用欠采样技术（如NearMiss）减少多数类样本，或在模型训练中结合过采样与欠采样（如SMOTEENN）。
	\item \textbf{算法层代价敏感学习：} 在训练模型（如逻辑回归、SVM、决策树）时，为少数类样本设置更高的误分类代价（如调整`class\_weight`参数），迫使模型更关注爆款样本的学习。
	\item \textbf{集成方法：} 使用EasyEnsemble或BalanceCascade等集成方法，通过构建多个平衡的子训练集来训练基分类器，然后集成结果。
\end{itemize}

\section{多模态特征融合}
小红书笔记包含文本、图像、视频、上下文等多种模态信息。如何有效融合这些异构特征是提升模型性能的关键。
\begin{itemize}
	\item \textbf{早期融合：} 将不同模态的特征在输入层进行拼接，形成一个大特征向量，然后输入到传统机器学习模型（如XGBoost）或深度学习模型（如MLP）中。这种方法简单，但可能忽略模态间的深层交互。
	\item \textbf{晚期融合：} 为每个模态单独训练一个子模型（如用BERT处理文本，用ResNet处理图像），然后将各子模型的预测概率进行加权平均或作为元特征输入到一个融合模型（如逻辑回归）中进行最终决策。
	\item \textbf{中级融合与注意力机制：} 设计深度学习架构，分别用不同的子网络处理不同模态的数据，在中间层进行特征交互与融合。引入跨模态注意力机制（Cross-modal Attention），让模型动态地决定在做出分类决策时，哪些模态的哪些部分信息更重要。例如，对于一篇美妆教程笔记，模型可能更关注文本中的步骤描述和图片中的效果对比。
\end{itemize}

\section{引入图神经网络建模社交传播}
笔记的热度不仅取决于其自身质量，还与发布者的社交网络和社区互动有关。虽然数据受限，但我们可以尝试构建简化的图结构。
\begin{itemize}
	\item \textbf{构建笔记-用户二部图：} 节点分为笔记和用户两类。边表示用户与笔记的交互（点赞、收藏、评论）。通过图神经网络（如GCN、GraphSAGE）学习节点嵌入，笔记节点的嵌入可以融合其内容特征和其在图中的传播结构信息（如被哪些高影响力用户互动过），从而提升分类精度。
	\item \textbf{构建用户相似图：} 基于用户兴趣或行为相似性构建用户之间的边。通过学习用户节点的嵌入，可以丰富发布者特征，更好地表征其粉丝群体的潜在偏好。
\end{itemize}

\section{利用预训练语言与视觉模型}
迁移学习能大幅提升文本和图像特征的表示能力。
\begin{itemize}
	\item \textbf{文本特征：} 使用在大规模语料上预训练的中文BERT、RoBERTa等模型，对笔记标题和正文进行编码，获取包含丰富语义信息的上下文向量表示，替代传统的TF-IDF或Word2Vec特征。
	\item \textbf{图像特征：} 使用在ImageNet等大型数据集上预训练的CNN模型（如EfficientNet, CLIP）提取图像的通用视觉特征。CLIP模型因其对齐了图像和文本语义空间，特别适合理解小红书笔记中图文匹配的程度。
\end{itemize}

\chapter{模型评估与比较}

为了科学评估和比较不同分类模型的性能，我们将采用一系列评估指标，并重点关注模型在“高热度”少数类上的识别能力。

\section{评估指标}
鉴于数据集的不平衡性，准确率（Accuracy）不是可靠的评估指标。我们将主要使用以下指标：
\begin{itemize}
	\item \textbf{精确率 (Precision)：} 模型预测为爆款的笔记中，真正是爆款的比例。高精确率意味着推送或推广的误判成本低。
	\item \textbf{召回率 (Recall)：} 所有真实爆款笔记中，被模型正确找出的比例。高召回率意味着不错过潜在的爆款机会。
	\item \textbf{F1-Score：} 精确率和召回率的调和平均数，是综合衡量模型在正类上性能的单一指标。我们主要参考宏平均F1（Macro-F1），它平等看待每一类，对少数类性能变化更敏感。
	\item \textbf{ROC曲线与AUC：} 绘制真正例率（TPR/Recall）随假正例率（FPR）变化的曲线。曲线下面积（AUC）衡量模型整体区分类别的能力，AUC越接近1，模型性能越好。ROC-AUC对类别不平衡相对不敏感。
	\item \textbf{PR曲线与AUC：} 绘制精确率随召回率变化的曲线。在不平衡数据中，PR曲线比ROC曲线更能反映模型在正类上的真实性能。PR曲线下面积（PR-AUC）越高，表明模型在保证高精确率的同时也能获得高召回率的能力越强。
	\item \textbf{混淆矩阵：} 直观展示模型在四个类别（TP, FP, FN, TN）上的具体表现，便于分析错误类型。
\end{itemize}

\section{模型比较与选择流程}
我们将遵循以下严谨流程进行模型比较与选择：
\begin{enumerate}
	\item \textbf{数据划分：} 将数据集按时间顺序划分为训练集、验证集和测试集（如7:2:1），以模拟现实中的时间外预测，避免未来信息泄露。
	\item \textbf{交叉验证：} 在训练集上使用时序交叉验证（TimeSeriesSplit）进行超参数调优和模型选择，确保评估的稳健性。
	\item \textbf{基准线：} 设定一个简单的基准线（如预测所有笔记为普通类，或基于粉丝数的简单规则模型），所有复杂模型性能需显著优于该基准线。
	\item \textbf{综合评估：} 在验证集上，综合比较各模型的宏平均F1、ROC-AUC、PR-AUC等指标。优先选择在PR-AUC和召回率上表现突出的模型，因为本研究更关注尽可能多地识别出潜在爆款（高召回），同时控制误判率（高精确率）。
	\item \textbf{测试集验证：} 最终选定1-3个最佳模型，在独立的测试集上进行最终性能报告，确保结果的可靠性。
	\item \textbf{可解释性分析：} 对最优模型（如XGBoost、带注意力机制的神经网络）进行可解释性分析，使用SHAP值、特征重要性图、LIME等方法，量化并可视化各特征对预测结果的贡献，从而得出可操作的、用于提升笔记热度的策略建议。
\end{enumerate}


\part{PartV: 小红书帖子热度预测模型}
\chapterimage{chapter_head_2.pdf} % Chapter heading image

\chapter{基准模型}

\section{预测问题描述}

小红书帖子热度预测的核心目标是预测帖子在未来某个时间点的观看量、点赞数、收藏数等互动指标，这是一个典型的回归预测问题。与简单的热门/非热门二分类预测不同，回归预测旨在提供连续的数值输出，这对于内容创作者优化发布策略、品牌方评估广告投放效果以及平台优化推荐算法具有更直接的指导意义。在社交媒体营销的实际应用中，热度预测的应用场景非常广泛。例如，在内容发布前，可以根据历史数据特征预测潜在热度，辅助创作者优化标题、标签和发布时间；在发布后短时间内，可以基于早期互动数据预测最终热度，为实时流量分配和推荐加权提供依据。

本研究的建模思路是，首先构建一个包含多维特征的输入向量，这些特征可以包括内容文本特征（标题情感、关键词密度、文本长度）、视觉特征（图片数量、图片质量评分、是否包含人脸/产品）、时间特征（发布时段、星期几、是否为节假日）、空间特征（发布地点热度、地理标签）、社交特征（粉丝数、历史发帖频率、账号垂直度）以及标签特征（tag数量、tag热度、tag与内容相关性）等。然后，选择合适的回归模型来学习这些特征与帖子热度指标之间的复杂非线性关系。我们将重点考察支持向量回归（SVR）、梯度提升树（GBDT）和深度学习模型（如LSTM、Transformer）在这一任务中的表现。SVR通过其独特的$\epsilon$-不敏感损失函数，能够在拟合精度和模型泛化能力之间取得良好的平衡，对社交媒体数据中的噪声和异常值具有较强的鲁棒性。而LSTM和Transformer则擅长捕捉时间序列数据中的长期依赖关系，能够从用户互动序列中学习到潜在的传播动态模式。通过对这些模型的构建、训练和评估，我们期望能够开发出预测精度高、稳健性强的小红书帖子热度预测模型，为后续的内容优化策略和精准营销提供坚实的基础。

为了全面评估所提出模型的性能，我们将构建一系列具有代表性的基准模型，涵盖从传统统计方法到现代机器学习和深度学习模型。这些基准模型将作为比较的基准，以衡量我们改进模型的有效性。

结合不同学习算法预测，建立一套多维度预测系统。

机器学习：KNN, SVR, Linear Regression, Ridge, Lasso, Elastic Net, Decision Tree, Extra Trees

深度学习：LSTM, GRU, BPNN, CNN, TextCNN, BERT, RoBERTa

集成学习：AdaBoost, GBDT, Bagging, Extra Trees, XGBoost, LightGBM, CatBoost, Random Forest

宽度学习：Broad Learning System (BLS)

\section{线性回归模型}

线性回归模型是回归分析中最基础、最经典的模型。它假设因变量（如帖子点赞数或观看量）与一个或多个自变量（如字数、图片数、tag数）之间存在线性关系。该模型的优点是形式简单、易于实现和解释，计算效率高，能够快速提供特征与热度之间的边际效应分析。在社交媒体分析领域，线性回归常被用作初步探索变量之间关系的工具，或作为更复杂模型的基准。例如，可以用来量化"每增加一张图片平均能带来多少额外点赞"或"晚间发布相比午间发布的热度差异"。

然而，其最大的局限性在于线性假设，这使其难以捕捉社交媒体数据中普遍存在的非线性关系（如tag数量与热度之间可能存在的倒U型关系）。此外，线性回归对异常值（如偶然爆红的 viral content）和多重共线性（如粉丝数与历史平均点赞数高度相关）较为敏感，这可能会影响其预测的准确性和稳健性。尽管如此，由于其简单性和可解释性，线性回归模型仍然是一个重要的基准，可以用来评估更复杂模型在处理非线性问题上的增益，并为特征重要性提供直观的理解。

\subsection{经典线性模型}
普通最小二乘法（OLS）线性回归将作为最基础的基准模型，用于建立特征与热度指标之间的线性映射关系。

\subsection{正则化线性模型}
引入L1正则化（Lasso）和L2正则化（Ridge）的线性模型，用于处理特征多重共线性问题并进行特征选择，防止在多维社交媒体特征场景下过拟合。

\section{支持向量回归（SVR）}

支持向量回归（Support Vector Regression, SVR）是支持向量机（SVM）在回归问题上的重要应用，它继承了SVM在处理小样本、非线性及高维模式识别问题中的诸多优势。SVR的核心思想是寻找一个最优的回归函数（或称为超平面），使得所有训练样本点与该函数之间的偏差最小，同时保证模型的泛化能力。其最大的创新在于引入了"$\epsilon$-不敏感损失函数"（$\epsilon$-insensitive loss function）。该函数设定了一个宽度为$2\epsilon$的"容忍带"，只有当预测值与真实值之间的误差超出这个容忍带时，才会计算损失并进行惩罚。这种机制使得模型对容忍带内的小误差不敏感，从而提高了对噪声数据的鲁棒性，并使得模型解具有稀疏性，即只有少数"支持向量"对最终的回归函数有贡献。

在小红书热度预测中，这种特性尤为重要，因为社交媒体数据往往包含大量噪声（如随机流量波动、机器人账号干扰）。SVR能够有效地捕捉内容特征（如情感极性、视觉美学）与热度指标之间复杂的非线性关系，克服了简单线性模型的局限性。

SVR模型有三个关键参数需要调优：
\begin{itemize}
	\item 容忍带宽度 ($\epsilon$)：$\epsilon$决定了对误差的容忍程度。$\epsilon$越小，模型对拟合精度的要求越高，可能导致模型复杂、过拟合；$\epsilon$越大，模型越简单，但可能欠拟合。在小红书数据中，由于热度分布通常呈现长尾特性，需要谨慎设置以避免对高热度异常值过度敏感。
	\item 惩罚系数 (C)：C控制了对超出容忍带的样本的惩罚力度。C越大，对误差的惩罚越重，模型倾向于将所有样本都拟合好，可能过拟合；C越小，允许更多的误差，模型泛化能力可能更强。考虑到小红书帖子热度的极端方差，C的选择需要平衡对 viral content 的拟合能力。
	\item 核函数 (Kernel Function)：核函数是SVR处理非线性问题的关键。通过核技巧，SVR可以将原始数据映射到更高维的特征空间，从而在新的空间中寻找线性回归函数。常用的核函数包括线性核、多项式核和高斯径向基函数（RBF）核。RBF核因其良好的性能和普适性，在处理文本、图像混合特征的社交媒体数据中被广泛使用。
\end{itemize}

实证研究表明，通过合理的参数调优和特征工程，SVR在预测社交媒体内容参与度方面能够取得优于传统线性模型的效果，特别是在中等规模数据集上表现稳定。

\section{集成学习模型（Random Forest, XGBoost, LightGBM）}

集成学习（Ensemble Learning）是一种通过构建并结合多个学习器来完成预测任务的机器学习方法。其核心思想是"三个臭皮匠，顶个诸葛亮"，通过集成多个弱学习器的预测结果，可以获得一个比任何单个学习器都更强大、更稳健的强学习器。在小红书热度预测领域，集成学习模型因其出色的性能、对高维特征的鲁棒性以及可解释性而备受青睐。

随机森林（Random Forest）是一种基于Bagging（Bootstrap Aggregating）思想的集成学习算法。它通过自助采样法（Bootstrap）从原始训练集中生成多个不同的子训练集，并在每个子训练集上独立地训练一棵决策树。在构建决策树的过程中，随机森林还引入了特征的随机选择，即每次分裂节点时，只从所有特征中随机选取一个子集来寻找最优分裂特征。这种双重随机性使得随机森林中的每棵树都具有差异性，从而降低了模型的方差，提高了泛化能力。随机森林对异常值和噪声不敏感，能够处理高维数据（如大量的tag特征、文本TF-IDF特征），并且可以评估特征的重要性，因此在理解哪些因素（如图片vs文本vs时间）对热度影响最大方面表现出色。

XGBoost（eXtreme Gradient Boosting）是一种基于Boosting思想的集成学习算法。与Bagging并行训练多个弱学习器不同，Boosting采用串行的方式，每一轮迭代都根据前一轮模型的预测结果来调整样本权重，使得新的模型能够更关注于之前被错误预测的样本（如那些"本该热门却表现平平"或"意外爆红"的帖子）。XGBoost是梯度提升决策树（Gradient Boosting Decision Tree, GBDT）的一种高效实现，它在GBDT的基础上进行了多项优化，包括使用二阶泰勒展开来近似损失函数、引入正则化项来控制模型复杂度、以及采用并行计算来加速训练过程。XGBoost以其卓越的预测性能、高效性和灵活性，在各种机器学习竞赛中屡获殊荣，也被广泛应用于社交媒体热度预测和广告效果评估中。

LightGBM（Light Gradient Boosting Machine）是另一种高效的梯度提升框架，采用基于直方图的决策树算法和叶子优先（Leaf-wise）的树生长策略，相比XGBoost在训练速度和内存效率上有显著优势，特别适合处理小红书这类大规模、高维度的社交媒体数据集。

\subsection{树集成模型一：Random Forest}
利用随机森林模型进行特征重要性分析，识别对观看量和点赞数影响最大的内容特征，为内容创作者提供可操作的优化建议。

\subsection{树集成模型二：XGBoost/LightGBM}
利用XGBoost和LightGBM构建高精度预测模型，通过网格搜索和贝叶斯优化进行超参数调优，捕捉特征间的复杂交互效应（如"高颜值图片+晚间发布+情感化标题"的组合效应）。

\section{深度学习模型（LSTM, TextCNN, BERT）}

深度学习模型，特别是循环神经网络（RNN）及其变体、卷积神经网络（CNN）以及预训练语言模型，在处理小红书多模态数据（文本+图像+时序）方面展现出强大的能力，因为它们能够自动从原始数据中学习复杂的模式和层次化的特征表示。

长短期记忆网络（Long Short-Term Memory, LSTM）是一种特殊的RNN，专门设计用于解决传统RNN在处理长序列时面临的梯度消失和梯度爆炸问题。LSTM通过引入一个精巧的门控机制（包括输入门、遗忘门和输出门）来控制信息流，使其能够有选择地记忆和遗忘信息，从而有效地捕捉时间序列中的长期依赖关系。在小红书热度预测中，LSTM可以用于建模用户互动的时间动态，例如，基于发布后1小时、2小时、6小时的点赞/收藏序列，预测24小时后的最终热度。这种早期热度序列往往包含了内容质量的重要信号。

卷积神经网络（Convolutional Neural Network, CNN）在图像识别领域取得了巨大成功，其强大的局部特征提取能力也被成功应用于文本分析（TextCNN）。在小红书场景下，TextCNN可以通过不同大小的卷积核捕捉标题和正文中的局部语义特征（如关键词组合、情感表达模式），有效识别"爆款标题"的文本模式。对于图像数据，可以使用预训练的CNN（如ResNet, VGG）提取视觉特征，评估图片质量、美学得分和内容相关性。

预训练语言模型（BERT, RoBERTa）通过大规模语料预训练，能够深度理解小红书文本内容的语义和上下文信息。利用BERT对标题和正文进行编码，可以获得高质量的文本表征，捕捉隐含的情感倾向、话题类别和写作风格，这些都是影响热度的重要因素。

此外，可以构建多模态融合模型（Multi-modal Fusion），将文本特征（BERT）、视觉特征（CNN）和结构化特征（元数据）在特征层或决策层进行融合，全面捕捉影响热度的多维度因素。

\subsection{时序深度学习模型：LSTM/GRU}
构建基于早期互动序列（发布后前N小时的点赞、评论、分享数）的LSTM模型，实现热度的早期预测和趋势外推。

\subsection{文本深度学习模型：BERT + TextCNN}
利用BERT提取深层语义特征，结合TextCNN捕捉局部关键词模式，构建高精度的文本内容质量评估和热度预测模型。

\section{其他模型}

\subsection{宽度学习系统（Broad Learning System, BLS）}
宽度学习是一种新兴的平坦神经网络架构，通过随机生成特征节点和增强节点，利用伪逆快速计算输出权重，具有训练速度极快、结构灵活的特点。在小红书实时热度预测场景中，BLS可以作为快速基准模型，支持模型的在线更新和增量学习。

\subsection{生存分析模型（Survival Analysis）}
考虑到部分小红书帖子可能长期持续获得流量（"长尾内容"），而另一些则快速衰减，可以引入生存分析中的Cox比例风险模型或AFT（Accelerated Failure Time）模型，不仅预测总热度，还预测热度的持续时间和衰减速度，为评估内容的长期价值提供依据。

\chapter{我们提出的改进模型}

在基准模型的基础上，我们将设计并实现一个或多个改进模型，以期在预测精度上取得突破。改进方向将主要集中在多模态数据融合、用户行为动态建模和上下文感知优化上。

\section{模型架构设计}

在基准模型的基础上，本研究将致力于提出改进模型，并探索如何将预测模型与内容创作策略进行有效融合，以期开发出具有更高传播效果的智能化内容优化系统。

改进模型方面，我们将从以下几个方向进行探索：

\begin{itemize}
	\item 多模态深度融合：现有的多模态方法通常采用简单的特征拼接或浅层融合。我们将探索更先进的融合机制，如基于注意力机制的多模态融合（Attention-based Fusion）、跨模态Transformer（Cross-modal Transformer），让模型自动学习文本、图像、时间、空间特征之间的复杂交互关系。例如，模型可以学习到"美食类内容在晚间配合高饱和度图片效果更佳"这样的隐性规则。
	
	\item 图神经网络（GNN）应用：构建用户-内容-标签的异构图网络，利用图神经网络捕捉社交关系、内容相似性和标签共现关系。通过消息传递机制，模型可以借鉴相似内容的历史表现，或利用用户的社交影响力来修正热度预测，解决冷启动问题。
	
	\item 多任务学习框架：不同于单一预测观看量或点赞数，我们将构建多任务学习模型，同时预测多个热度指标（观看量、点赞数、收藏数、评论数、分享数）。通过共享表征层，模型可以学习到不同互动行为之间的内在关联（如"高收藏内容通常评论较少但转化率较高"），提高整体预测精度。
	
	\item 时序卷积网络（TCN）与Transformer：引入Temporal Convolutional Network和Transformer架构替代LSTM，利用其并行计算优势和长距离依赖捕捉能力，更有效地建模热度传播的时间动态，特别是捕捉突发流量和病毒式传播模式。
	
	\item 因果推断集成：结合因果推断方法（如双重机器学习、倾向得分匹配），不仅预测相关性，还尝试识别因果关系，区分"因为发布了某类内容所以热门"和"因为账号本身有影响力所以热门"，为内容策略提供更 robust 的建议。
\end{itemize}

策略融合方面，我们将探索如何将回归预测模型的输出直接转化为内容优化策略。例如：

\begin{itemize}
	\item 实时发布决策支持：构建一个策略系统，当模型预测某篇草稿在特定时间、搭配特定标签发布能获得更高热度时，向创作者推荐最佳发布时机和标签组合。
	
	\item A/B测试优化：利用预测模型进行虚拟A/B测试，在内容发布前预测不同标题、封面图版本的表现，辅助创作者选择最优方案。
	
	\item 动态流量分配：对于品牌方和MCN机构，预测模型可以用于评估不同KOL的内容投放效果，优化广告预算分配，识别潜在的高性价比投放对象。
	
	\item 内容冷启动策略：针对新账号或新内容，利用预测模型识别"潜力内容"，建议平台给予初始流量扶持，提高优质内容的发现效率。
\end{itemize}

通过这种深度融合，我们期望能够开发出比传统内容运营更智能、更数据驱动的小红书营销优化系统，实现从"经验驱动"到"算法驱动"的转变。


\part{Part VI: 实验设计与结果分析}
\chapterimage{chapter_head_2.pdf} % Chapter heading image

\chapter{实验设置与评估指标}

\section{实验环境与超参数设置}
本研究所有实验均在统一的计算环境下进行，以确保结果的可比性。模型训练使用Python作为主要编程语言，并依赖Scikit-learn、XGBoost、LightGBM、PyTorch等机器学习框架。为公平比较，所有模型均在相同的数据划分（训练集：验证集：测试集 = 70\%:15\%:15\%）上进行训练与评估，且采用相同的随机种子以确保可复现性。

各模型的关键超参数均通过网格搜索（Grid Search）或随机搜索（Random Search）结合时序交叉验证进行优化。主要优化目标为在验证集上最小化预测“点赞数”的均方根误差。例如，XGBoost模型重点调节了学习率、最大深度、子样本比例等；LSTM模型则优化了隐藏层维度、学习率、Dropout率等。

\section{回归预测评估指标体系}
为全面、客观地评估各模型对“观看量”和“点赞数”的预测性能，我们采用了一套多维度的评估指标。这些指标从不同角度衡量预测误差，以克服单一指标的局限性。

\begin{table}[H]
	\centering
	\caption{回归预测任务主要评估指标}
	\begin{tabularx}{\textwidth}{l l X X X}
		\toprule
		\textbf{评估指标} & \textbf{缩写} & \textbf{定义} & \textbf{优点} & \textbf{局限性} \\
		\midrule
		平均绝对误差 & MAE & $\frac{1}{n}\sum_{i=1}^{n} |y_i - \hat{y}_i|$ & 直观，解释性强；对异常值不敏感 & 无法反映误差的相对大小 \\
		均方根误差 & RMSE & $\sqrt{\frac{1}{n}\sum_{i=1}^{n}(y_i - \hat{y}_i)^2}$ & 对大误差惩罚更重；数学性质优良 & 对异常值敏感；数值尺度依赖原始数据 \\
		平均绝对百分比误差 & MAPE & $\frac{1}{n}\sum_{i=1}^{n} \left|\frac{y_i - \hat{y}_i}{y_i}\right| \times 100\%$ & 无量纲，便于不同量纲任务比较 & 当$y_i$接近或等于0时，指标会趋向无穷大或不稳定 \\
		对称平均绝对百分比误差 & sMAPE & $\frac{200\%}{n}\sum_{i=1}^{n} \frac{|y_i - \hat{y}_i|}{|y_i| + |\hat{y}_i|}$ & 克服了MAPE在零值附近的问题，定义更对称 & 当预测值和真实值均为零时，分母为零 \\
		决定系数 & $R^2$ & $1 - \frac{\sum_{i=1}^{n}(y_i - \hat{y}_i)^2}{\sum_{i=1}^{n}(y_i - \bar{y})^2}$ & 衡量模型对数据方差的解释能力；值域[0,1]，越接近1越好 & 对模型复杂度不敏感，可能随特征增加而虚高 \\
		\bottomrule
	\end{tabularx}
	\label{tab:metrics}
\end{table}

\noindent \textbf{指标选择说明：}
\begin{itemize}
	\item \textbf{MAE和RMSE}：作为核心的绝对误差指标，用于评估预测的精确度。RMSE对较大预测误差更为敏感，能反映模型的稳健性。
	\item \textbf{MAPE}：由于“观看量”和“点赞数”数值范围大，MAPE有助于评估相对误差，更贴近业务中对“预测偏差百分比”的直观理解。
	\item \textbf{$R^2$}：用于评估模型相对于简单基准（如使用均值预测）的优越性，衡量模型的整体拟合优度。
\end{itemize}
所有指标的最终报告值均基于模型在独立测试集上的表现。

\chapter{基准模型预测结果}

\section{基准模型性能比较}
我们将所有基准模型在测试集上对“观看量”（取对数后）和“点赞数”（取对数后）的预测结果进行了系统评估。表\ref{tab:baseline_results}汇总了主要性能指标。

\begin{table}[H]
	\centering
	\caption{基准模型预测性能对比（测试集）}
	\begin{tabular}{l c c c c c c}
		\toprule
		
		\bottomrule
	\end{tabular}
	\label{tab:baseline_results}
\end{table}

\noindent \textbf{结果分析：}
\begin{enumerate}
	\item \textbf{树模型优势显著}：XGBoost和随机森林在各项指标上全面领先于线性模型和SVR，这证实了小红书笔记热度与多维特征之间存在复杂的非线性关系和交互作用，而树模型擅长捕捉此类模式。
	\item \textbf{深度学习模型表现}：MLP表现与随机森林相近，但略逊于XGBoost。纯LSTM模型（仅使用发布者历史互动序列等时序特征）表现中等，说明在静态回归预测任务中，丰富的结构化特征比纯序列特征信息量更大。
	\item \textbf{线性模型仍有价值}：线性回归和SVR虽不及树模型，但$R^2$仍达到0.44左右，远优于基线，表明特征与目标间存在较强的线性可解释成分。其简洁性和高可解释性在初步分析中仍有应用价值。
\end{enumerate}

\section{预测结果可视化分析}
为更直观地展示模型预测效果，我们选取了在“点赞数”预测上表现最佳的XGBoost模型进行分析。

%\begin{figure}[H]
%	\centering
%	\includegraphics[width=0.8\textwidth]{figures/prediction_scatter.png}
%	\caption{XGBoost模型预测点赞数（对数） vs. 真实点赞数（对数）散点图}
%	\label{fig:scatter}
%\end{figure}

从模型预测值与真实值的对比结果可观察到：
\begin{itemize}
	\item 在中等热度区域（log(点赞数)在3-6之间），预测点紧密分布在对角线附近，说明模型预测较为准确。
	\item 在极低热度区域（长尾部分），预测存在一定高估倾向；在极高热度区域（爆款），预测存在一定低估倾向。这是由于极端样本数量稀少，模型难以完全学习其复杂模式，也是热度预测领域的常见挑战。
\end{itemize}

\chapter{改进模型实验与消融分析}

\section{多任务学习与多模态融合模型结果}
我们提出的改进模型（多模态融合 + 多任务学习框架）在测试集上取得了最佳性能，具体结果如表\ref{tab:improved_results}所示。

\begin{table}[H]
	\centering
	\caption{改进模型与最佳基准模型性能对比}
	\begin{tabular}{l c c c c}
		\toprule
		
		\bottomrule
	\end{tabular}
	\label{tab:improved_results}
\end{table}

\noindent \textbf{关键发现：}
\begin{enumerate}
	\item \textbf{多模态融合有效}：即使是在单一任务下，改进模型（融合了深度文本、图像特征及图网络嵌入）的$R^2$也显著优于仅使用传统特征工程的XGBoost，证明深度多模态特征提供了额外的预测信息。
	\item \textbf{多任务学习带来增益}：采用多任务学习框架后，两个目标的预测$R^2$进一步同步提升。这表明“观看量”和“点赞数”这两个任务共享底层信息，联合学习能使模型学习到更通用的内容质量表示，从而提高了泛化能力。
\end{enumerate}

\section{消融实验分析}
为验证改进模型中各核心组件的有效性，我们设计了系统的消融实验，结果如表\ref{tab:ablation}所示（以多任务框架下“点赞数”预测$R^2$为例）。

\begin{table}[H]
	\centering
	\caption{改进模型消融实验分析}
	\begin{tabular}{lc}
		\toprule
		\textbf{模型变体} & \textbf{点赞数预测 $R^2$ $\uparrow$} \\
		\midrule
		\textbf{完整模型} & \textbf{0.721} \\
		- 移除跨模态注意力模块 & 0.708 \\
		- 移除GNN社交传播特征 & 0.714 \\
		- 使用TF-IDF替代BERT文本特征 & 0.698 \\
		- 移除多任务学习（改为独立训练） & 0.709 \\
		- 所有组件均移除（相当于MLP） & 0.631 \\
		\bottomrule
	\end{tabular}
	\label{tab:ablation}
\end{table}

\noindent \textbf{组件贡献度解读：}
\begin{itemize}
	\item \textbf{跨模态注意力贡献最大}：移除后性能下降最明显（0.013），说明该模块有效促进了图文信息的深度融合，对理解内容整体吸引力至关重要。
	\item \textbf{GNN特征与多任务学习次之}：各自移除导致约0.007的性能下降。GNN特征编码了社交传播潜力，多任务学习促进了通用表示学习，二者均为有效改进。
	\item \textbf{深度文本特征重要}：用TF-IDF替代预训练的BERT文本特征导致性能显著降低，证实了深度语义理解对内容质量评估的价值。
\end{itemize}

\chapter{特征重要性分析与策略洞见}

\section{全局特征重要性排名}
基于XGBoost模型（因其优秀的特征重要性评估功能），我们得出了影响小红书笔记热度的全局特征重要性排名（综合观看与点赞预测）。

%\begin{figure}[H]
%	\centering
%	\includegraphics[width=0.9\textwidth]{figures/feature_importance.png}
%	\caption{基于XGBoost模型的特征重要性排名（Top 15）}
%	\label{fig:feature_importance}
%\end{figure}

\noindent \textbf{核心发现：}
\begin{enumerate}
	\item \textbf{发布者影响力是基石}：“发布者粉丝数（对数）”和“发布者历史平均点赞数”高居前两位，这印证了初始流量池的至关重要性。
	\item \textbf{内容质量与形式并重}：“标题情感积极度”、“正文可读性分数”、“封面图美学评分”和“是否包含视频”均位列前十，说明优质内容本身是获得互动的根本。
	\item \textbf{发布策略影响显著}：“发布时段（晚间）”、“是否含热门标签”排名靠前，表明优化运营策略能有效放大内容效果。
	\item \textbf{视觉先行}：图片相关特征（“封面图清晰度”、“图片数量”）的重要性普遍高于纯文本长度特征，凸显了小红书的“视觉种草”平台特性。
\end{enumerate}

\section{局部解释与案例研究}
我们使用SHAP值对具体笔记的预测结果进行局部解释。以一篇实际预测效果较好的美妆教程笔记为例：
\begin{itemize}
	\item \textbf{正向驱动因素（拉升预测值）}：1) “发布者粉丝数众多”（高SHAP值）；2) “封面图为精心设计的对比图”；3) “文案使用了‘保姆级’等强吸引力词汇”。
	\item \textbf{负向驱动因素（拉低预测值）}：1) “发布在周三上午的流量低谷期”；2) “未使用任何当期平台热门标签”。
\end{itemize}
这种局部解释能为内容创作者提供直接、可操作的优化建议，例如：即使内容优质，也需注意发布时间和话题借势。

\chapter{统计显著性检验与模型排名}

\section{Friedman检验与Nemenyi事后检验}
为判断不同模型之间的性能差异是否具有统计学显著性，我们基于各模型在5折交叉验证中获得的“点赞数预测$R^2$”排名，进行了Friedman检验。
\begin{itemize}
	\item \textbf{Friedman检验结果}：计算得到的统计量显著大于临界值（$p < 0.01$），表明在95\%的置信水平下，可以拒绝“所有模型性能相同”的原假设，即至少有一个模型与其他模型存在显著差异。
	\item \textbf{Nemenyi事后检验}：随后进行的Nemenyi检验用于两两比较。其生成的临界差异（CD）图（此处以文字描述）显示：
	\begin{enumerate}
		\item 改进模型与XGBoost、随机森林等基准模型之间的差异超出了临界差异，性能提升具有统计显著性。
		\item XGBoost、随机森林、MLP三者处于同一个性能分组，它们之间的差异不显著。
		\item 线性回归、SVR等模型处于另一个显著较低的性能分组。
	\end{enumerate}
\end{itemize}

\section{实验结论总结}
综合以上实验结果，可得出以下主要结论：
\begin{enumerate}
	\item \textbf{模型性能层面}：融合多模态特征与多任务学习的深度学习改进模型，在预测小红书笔记“观看量”和“点赞数”上达到了最佳性能，且其优势具有统计显著性。
	\item \textbf{影响因素层面}：笔记热度是发布者影响力、内容质量（尤其是视觉质量）、运营策略（时间、标签）三者共同作用的结果。其中，发布者基础决定了起点，内容质量决定了转化深度，而运营策略则影响了被推荐的广度。
	\item \textbf{实践指导层面}：
	\begin{itemize}
		\item 对于平台新账号，首要任务是提升粉丝基数，或寻求与KOL合作以利用其流量池。
		\item 内容创作应坚持“视觉优先”，在封面图和内容图片上投入精力。
		\item 发布前务必进行“热点校准”，通过使用热门标签、选择黄金发布时间来最大化内容的初始曝光概率。
	\end{itemize}
\end{enumerate}
本研究的模型与结论，为量化评估内容价值、数据驱动内容营销提供了有效的工具和明确的优化方向。






\part{Part VI: 结论与展望}
\chapterimage{chapter_head_2.pdf} % Chapter heading image





% Bibliography and index output are intentionally omitted until the report
% contains actual \cite and \index commands.  Running BibTeX/MakeIndex on the
% empty template sections caused latexmk to fail despite a valid main PDF.

%----------------------------------------------------------------------------------------

\end{document}
